% main_review2.tex: main (18).tex with the second round of review comments answered.
% Red: the comments, kept. Blue: their answers and this round's other edits. Removed text
% does not show.
% ferry_paper.tex: GENERATED by DeveloperDocs/paper/assemble.py from front.tex,
% methods.tex, setup.tex, evaluation.tex, back.tex and results/exp5/paper/*.tex.
% Edit those sources and rerun the script; edits here are overwritten.
% =============================================================================
% Front matter of the FeRRy paper: preamble, title, abstract, introduction,
% Table 1 and related works. Extracted on 8 Oct 2026 from Hermes_paper (6).tex
% (the live text only; commented and \iffalse'd drafts dropped), with a revised
% introduction. DeveloperDocs/paper/assemble.py appends methods.tex,
% evaluation.tex and back.tex to this file to build ferry_paper.tex.
% =============================================================================
\documentclass[conference]{IEEEtran}
\IEEEoverridecommandlockouts

% Text, Language, and URLs
\usepackage[utf8]{inputenc}
\usepackage[english]{babel}
\usepackage[hyphens]{url}
\usepackage{cite}

% Math and Symbols
\usepackage{amsmath,amssymb,amsfonts} % amssymb is included here
\usepackage{textcomp}

% Figures, Subfigures, and Colors
\usepackage{float}
\usepackage{graphicx}
\usepackage{subcaption}
\usepackage[table]{xcolor} % Updated to include [table] option

% Tables and Formatting
\usepackage{multirow}
\usepackage{tabularx}
\usepackage{booktabs} % Already existed
\usepackage{makecell} % Already existed
\usepackage{diagbox}
\usepackage{multicol}
\usepackage{fancyhdr}
\usepackage{colortbl} % Added for cell coloring
\usepackage{array}    % Added for advanced array/table options

% Custom Table Columns
\newcolumntype{L}{>{\raggedright\arraybackslash}X}
\newcolumntype{P}[1]{>{\centering\arraybackslash}m{#1}}

% Algorithms (Using algorithm + algpseudocode exclusively to avoid conflicts)
\usepackage{algorithm}
\usepackage{algpseudocode}

% TikZ and PGFPlots
\usepackage{tikz}
\usetikzlibrary{positioning, arrows.meta, calc}
\usepackage{pgfplots}
\pgfplotsset{compat=1.18} 

% Custom command definitions (defined exactly once)
\newcommand{\threat}[1]{\text{#1}}
\newcommand{\securitymechanism}[1]{\mathcal{S}_{\text{#1}}}
\newcommand{\protocol}[1]{\mathcal{C}_{\text{#1}}}
\newcommand{\experiment}[1]{\mathcal{E}_{\text{#1}}}
\newcommand{\performance}[1]{P_{\text{#1}}}

% FeRRy macros (methods and evaluation)
\newcommand{\sysname}{FeRRy}
\newcommand{\bbar}{\bar{b}}
\newcommand{\Tnom}{T_{\mathrm{nom}}}

\pagestyle{plain}

% \def\BibTeX{{\rm B\kern-.05em{\sc i\kern-.025em b}\kern-.08em
%     T\kern-.1667em\lower.7ex\hbox{E}\kern-.125emX}}


\makeatletter
\newcommand{\linebreakand}{%
  \end{@IEEEauthorhalign}
  \hfill\mbox{}\par
  \mbox{}\hfill\begin{@IEEEauthorhalign}
}
\makeatother

\begin{document}
% TODO(title): the paper now presents FeRRy; the old title named HERMES.
\title{FeRRy: RF-Aware Mission Planning for Federated Learning with Drones as Data Mules}
% Authors: omitted for double-anonymous review. The \author block (names,
% affiliations, emails) is kept in the Overleaf copy, not in this public repo;
% restore it there for the camera-ready.

\maketitle

% Abstract: the short hybrid of abstract_revision.md, adopted 8 Oct 2026
% (the document keeps the earlier versions and the audit of the original).
\begin{abstract}
Federated learning (FL) keeps training data on local devices, but it still
requires those devices to exchange model updates with a server. In disaster,
remote, and disrupted environments, some devices may be out of range or have
unreliable connections, causing their updates to be missed. UAVs can act as data
mules and physically carry model updates between these devices and the server,
but each UAV has limited mission time and its communication range affects where
it must fly. We present FeRRy, a Federated RF-aware Routing framework that selects the radio
band and UAV route for each mission. Before flight, FeRRy chooses a communication
range and route based on device deadlines, expected communication time, and the
mission budget. During flight, it can change the radio band or revise the
remaining route when the observed channel differs from the planned condition.
The UAV collects device updates, carries them back to the edge server, and later
delivers the updated global model. We evaluate FeRRy using a multi-process prototype that trains a real intrusion
detection model on CICIoT2023. The server, UAVs, and devices run as separate TCP
processes, while UAV movement, wireless conditions, mission time, and energy are
simulated. We compare FeRRy with seven scheduling approaches, including FedEx,
Oort, FedCS, and a Whittle index scheduler based on update age. With six devices,
one UAV, a 1~MB model update, and a moderate mission budget, FeRRy reaches the
target accuracy 23--32\% sooner than the baselines. The main gain comes from
selecting a narrow band with longer range, which greatly reduces UAV travel
between devices. \textcolor{blue}{Those baselines use the wide band; pinned to the narrow band, FedEx and a
deadline-based baseline are not significantly slower than FeRRy.} Under a tight budget, FeRRy also keeps updates fresher and stays
within the mission budget in 71\% of missions; the remaining missions exceed the
budget by 33~s on average. The advantage becomes smaller as model updates grow or
more UAVs are added: at 10~MB FeRRy selects the wide band and its lead is no
longer statistically significant, while with two or three UAVs most approaches
perform similarly.
\end{abstract}

\begin{IEEEkeywords}
Decentralized learning, federated learning, UAV data mules, tactical edge computing, mobile relay networks
\end{IEEEkeywords}

\section{Introduction}
\label{sec:introduction}

Disaster response, wildfire monitoring, agricultural sensing and tactical
operations increasingly rely on edge devices deployed where fixed infrastructure
is sparse, damaged or absent~\cite{zeng2016wireless, song2022comprehensive}.
These are also the settings most prone to denied, disrupted, intermittent and
limited (DDIL) connectivity~\cite{gupta2015survey}. Federated learning (FL) suits
them, because devices exchange model updates rather than raw
data~\cite{mcmahan2017communication, kostage2025hifins}, but it still needs every
device to reach the server: when a device cannot, its update is lost, and rounds
close short, time out, or fail to produce a usable model~\cite{li2020federated}.
The gap is not faster FL communication but federation that continues when the
network itself is fragmented.

UAVs acting as data mules can bridge disconnected regions by carrying model
updates between devices and servers~\cite{zhang2025latency, arsalaan2024uavs}.
However, a mule has limited flight time and energy, and a contact is useful only
if an update is ready and the link can complete the transfer when the mule
arrives. Hierarchical FL, UAV-assisted FL, physical model transport, RF-aware
networking, and client selection each address part of this
problem~\cite{liu2020client, lin2023delay, parasnis2023connectivity,
yemini2023robust, mestoukirdi2022uav, bian2025indirect, wang2018deep,
cui2019multi, nishio2019client, lai2021oort}, but they generally assume
persistent reachability, fixed aggregation points, or communication choices
that are separate from the UAV route.

We present \sysname{}, a Federated RF-aware Routing framework whose central
observation is that the mule's communication range can be part of the mission
decision. A narrower communication band reaches farther but transfers more
slowly. It may allow the UAV to serve several devices from one location, while
a wider band transfers faster but may require the UAV to fly closer to each
device. The better choice depends on how many devices need service, where they
are located, and how much mission time remains. \sysname{} therefore selects
the communication band as part of the route planning before each flight.

\sysname{} makes decisions before and during the mission. Before takeoff, it
searches the available band classes and routes under the same feasibility test
for device deadlines and the mission budget. It favors plans that serve devices
that have waited longer for service. During flight, the mule uses the observed
channel to select a faster band when possible, reconsider the next stop, and
revise the remaining route when the original plan no longer fits within the
mission budget. Updates collected by each mule are combined before returning to
the server, and updates from multiple mules are then incorporated into the
global model. We evaluate three variants: FeRRy-Plan follows the route selected
before takeoff, FeRRy-Adapt adds deterministic in-mission adaptation, and
FeRRy-RL uses a learned rule for the in-mission decision.

In a multi-process prototype that trains a real intrusion-detection model on
CICIoT2023~\cite{CiCIoT2023}, with mission time, wireless conditions, and energy
simulated, we compare \sysname{} with seven scheduling approaches under an
analysis plan specified before each study. With six devices, one mule, and a
1~MB update, \sysname{} reaches the target accuracy 23--32\% sooner than the
baselines because its plan selects a narrow band with longer range and reduces
the travel between stops from about 140\,s to under 20\,s. When restricted to
the wide band, its performance is similar to FedEx. \textcolor{blue}{Conversely, FedEx and a deadline-based baseline pinned to the narrow band match
\sysname{} restricted to that band, so the gain in this setting comes from the RF
class rather than from route planning.} Under tight budgets,
\sysname{} also keeps updates fresher than age-based schedulers. The learned
in-mission rule does not improve over the deterministic rule, showing that most
of the gain comes from the mission plan itself. We also identify where the
advantage becomes smaller. The planner exceeds the tight mission budget in
29\% of missions, does not consider device training time, and provides less
benefit as update size or the number of mules increases. At 10~MB, it selects
the wide band and its lead is no longer statistically significant, while with
two or three mules most approaches perform similarly.

The contributions of this paper are:

\begin{itemize}

\item \textbf{Communication range as a mission decision.}
We treat the communication band, which determines both contact range and
transfer rate, as part of the route-planning decision under a finite mission
budget.

\item \textbf{Mission planning with in-flight adaptation.}
\sysname{} uses the same feasibility test before and during flight to account
for device deadlines, participant coverage, communication time, and the
remaining mission budget. We evaluate both deterministic and learned in-mission
decisions.

\item \textbf{A prototype and simulator with shared mission logic.}
Our multi-process prototype trains and exchanges a real model, while FerrySim
uses the same planner, feasibility test, RF model, and mission logic for learned
policies and larger deployments of up to 96 devices.

\item \textbf{A controlled comparison with existing approaches.}
We compare \sysname{} with seven scheduling approaches using paired random
seeds, confidence intervals, and multiplicity-corrected statistical tests. We
report where \sysname{} improves performance, which parts of the design
contribute to those gains, and where its advantage becomes smaller.


\end{itemize}

Section~\ref{sec:related_works} reviews related work,
Section~\ref{sec:method} presents \sysname{}'s design,
Section~\ref{sec:evaluation} the experimental setup and design,
Section~\ref{sec:results} the results, and Section~\ref{sec:conclusion}
concludes.

\definecolor{goodgreen}{RGB}{178,210,150}
\definecolor{ferrygreen}{RGB}{105,160,90}
\definecolor{badred}{RGB}{230,170,160}
\definecolor{nayellow}{RGB}{240,210,120}

\newcommand{\cmarkc}{\cellcolor{goodgreen}\textcolor{black!55}{\scriptsize$\checkmark$}}
\newcommand{\xmarkc}{\cellcolor{badred}\textcolor{black!40}{\scriptsize$\times$}}
\newcommand{\tmarkc}{\cellcolor{nayellow}\textcolor{black!65}{\scriptsize$\triangle$}}
\newcommand{\cmarkF}{\cellcolor{ferrygreen}\textcolor{white}{\textbf{\scriptsize$\checkmark$}}}

\begin{table*}[!t]
\caption{\small \textsc{Comparison of Related FL, UAV Mobility, and RF-Aware Coordination Approaches}}
\label{tab:ferry-comparison}
\centering
\resizebox{\textwidth}{!}{%
\begin{tabular}{l c c c c c c}
\toprule
\multicolumn{1}{c}{Framework}
&
\begin{tabular}[c]{@{}c@{}}Physical Model\\ Transport\end{tabular}
&
\begin{tabular}[c]{@{}c@{}}Explicit RF\\ Decision\end{tabular}
&
\begin{tabular}[c]{@{}c@{}}Mobility / Route\\ Decision\end{tabular}
&
\begin{tabular}[c]{@{}c@{}}Deadline /\\ Freshness Aware\end{tabular}
&
\begin{tabular}[c]{@{}c@{}}In-Mission\\ Replanning\end{tabular}
&
\begin{tabular}[c]{@{}c@{}}Joint FL--Mobility/RF\\ Coordination\end{tabular}
\\
\midrule

Hierarchical FL
{\scriptsize \cite{liu2020client}}
& \xmarkc & \xmarkc & \xmarkc & \xmarkc & \xmarkc & \xmarkc \\

Delay-Aware HFL
{\scriptsize \cite{lin2023delay}}
& \xmarkc & \xmarkc & \xmarkc & \tmarkc & \xmarkc & \xmarkc \\

FedCS
{\scriptsize \cite{nishio2019client}}
& \xmarkc & \xmarkc & \xmarkc & \cmarkc & \xmarkc & \xmarkc \\

Oort
{\scriptsize \cite{lai2021oort}}
& \xmarkc & \xmarkc & \xmarkc & \cmarkc & \xmarkc & \xmarkc \\

UAV Multi-Community FL
{\scriptsize \cite{mestoukirdi2022uav}}
& \xmarkc & \xmarkc & \cmarkc & \tmarkc & \xmarkc & \cmarkc \\

FedEx
{\scriptsize \cite{bian2025indirect}}
& \cmarkc & \xmarkc & \cmarkc & \tmarkc & \xmarkc & \cmarkc \\

Model-Aided Multi-UAV RL
{\scriptsize \cite{chen2023model}}
& \xmarkc & \tmarkc & \cmarkc & \xmarkc & \cmarkc & \xmarkc \\

Deadline-Sensitive Multi-UAV
{\scriptsize \cite{mondal2025deep}}
& \xmarkc & \tmarkc & \cmarkc & \cmarkc & \cmarkc & \xmarkc \\

DRL Multichannel Access
{\scriptsize \cite{wang2018deep}}
& \xmarkc & \cmarkc & \xmarkc & \xmarkc & \xmarkc & \xmarkc \\

MARL UAV Resource Allocation
{\scriptsize \cite{cui2019multi}}
& \xmarkc & \cmarkc & \tmarkc & \xmarkc & \xmarkc & \xmarkc \\

Value-Aware UAV FL
{\scriptsize \cite{cui2023data}}
& \xmarkc & \xmarkc & \xmarkc & \cmarkc & \xmarkc & \xmarkc \\

UAV-Enabled Asynchronous FL
{\scriptsize \cite{zhai2024uav}}
& \xmarkc & \tmarkc & \cmarkc & \cmarkc & \xmarkc & \cmarkc \\

FedBuff
{\scriptsize \cite{nguyen2022federated}}
& \xmarkc & \xmarkc & \xmarkc & \cmarkc & \xmarkc & \xmarkc \\

Async-HFL
{\scriptsize \cite{yu2023async}}
& \xmarkc & \xmarkc & \xmarkc & \cmarkc & \xmarkc & \xmarkc \\

\midrule

\textbf{FeRRy (ours)}
& \cmarkF
& \cmarkF
& \cmarkF
& \cmarkF
& \cmarkF
& \cmarkF \\

\bottomrule
\end{tabular}%
}

\vspace{2pt}
{\footnotesize
\cmarkc\ Present
\quad
\xmarkc\ Absent
\quad
\tmarkc\ Partial / indirect
}
\end{table*}


% Architecture figure: OUT of the build (8 Oct 2026). The drawing,
% HERMES_Case_Scenarios_compact.drawio, is the HERMES design: "WITH HERMES"
% labels, a DDQN choosing among carrier frequencies (FeRRy's band classes share
% one carrier and differ in bandwidth; nothing in FeRRy's RF layer is learned),
% and the old gated job scheduler. Redraw it to match Section III (plan clock,
% flight clock, merges), then uncomment this block and restore
% "A cluster (Fig.~\ref{fig:architecture})" in methods.tex's Setting paragraph.
% \begin{figure*}[t]
%     \centering
%     \includegraphics[trim=0cm 0cm 0cm 0cm, clip, scale=0.35]{Figures/HERMES_Case_Scenarios_compact.drawio.png}
%     \caption{\sysname{} overview: field devices, UAV data mules, and the edge server.}
%     \label{fig:architecture}
% \end{figure*}

\section{Related Works}
\label{sec:related_works}

Federated learning reduces centralized data movement by allowing participants to train locally and exchange model updates rather than raw datasets. FedAvg established this model-averaging paradigm, while FedProx improved robustness to statistical and systems heterogeneity \cite{mcmahan2017communication,li2020federated}. Hierarchical approaches subsequently introduced intermediate aggregation to reduce communication cost and accommodate heterogeneous edge environments. Aggregation-oriented hierarchical FL performs partial aggregation at edge servers before cloud synchronization, while delay-aware hierarchical FL explicitly accounts for communication delay between edge and cloud tiers \cite{liu2020client,lin2023delay}. Connectivity-aware semi-decentralized FL further exploits time-varying device-to-device links, and collaborative relaying allows well-connected clients to forward updates for clients experiencing intermittent server connectivity \cite{parasnis2023connectivity,yemini2023robust}. These approaches improve FL under unreliable communication, but they still rely on a communication graph in which model state is eventually forwarded through wireless links. FeRRy targets the more disruptive case in which such an end-to-end path may not exist and model state must instead be physically transported between isolated participants.

UAV-assisted FL introduces mobility into the learning infrastructure, but the role assigned to the UAV varies considerably. A UAV may operate as a mobile FL server whose routing is optimized with bandwidth, computation, and transmit power \cite{zhang2025latency}, or as a mobile orchestrator that visits discrete locations while optimizing routing and device scheduling altogether across multiple FL communities \cite{mestoukirdi2022uav}.  These systems use mobility to restore or improve communication between ground participants and the learning infrastructure. The closest architectural precedent to FeRRy is FedEx, which removes direct client-server connectivity and uses mobile transporters to physically carry model information \cite{bian2025indirect}. FedEx establishes physical model transport and relates FL performance to transporter travel time, but assumes preplanned tours and fixed communication conditions without adapting service to changing RF conditions, per-device deadlines, or in-mission contact opportunities. FeRRy extends this setting by treating reachability, service time, participant urgency, and communication conditions as part of the mission decision itself.

Another research direction studies UAV trajectory planning together with communication-aware data collection. Earlier formulations have demonstrated that trajectory and communication should not be optimized independently, and should jointly optimize multi-UAV trajectories, user scheduling, association, and transmit power for wireless service \cite{wu2018joint}, while more recent learning-based approaches move these decisions online. Multi-UAV controllers have used SNR, reachability, remaining device data, relative distance, and battery state to learn trajectories \cite{chen2023model}, and subsequent work has learned 3-D trajectories and device scheduling \cite{gao2024federated}, incorporated collection deadlines, link quality, and energy constraints through MAPPO \cite{mondal2025deep}, further optimized trajectory and task offloading under AoI, latency, and energy objectives \cite{wang2025joint}, and considered battery and interference constraints in federated-RL data collection \cite{kumar2026battery}. RF adaptation has likewise been studied through deep reinforcement learning for dynamic multichannel access \cite{wang2018deep} and multi-agent resource allocation over users, power levels, and subchannels in UAV networks \cite{cui2019multi}. These studies establish that communication-aware, and energy-aware multi-UAV trajectory learning is already well studied. FeRRy therefore does not claim RF-aware trajectory planning itself as novel, instead, its distinction is that the communication choice directly changes the feasibility and duration of transporting FL updates, while participant urgency determines which contact is worth serving next.

Client-selection and freshness-aware scheduling provides another part of this coordination problem. Existing methods select participants according to round deadlines and resource availability \cite{nishio2019client}, favor clients with higher training loss \cite{cho2022towards}, or use statistical utility and participation history observed in previous rounds \cite{lai2021oort}. This distinction between pre-selection and historical information is particularly important for physical data muling because isolated participants cannot continuously report their current state before the UAV reaches them. Related UAV-FL work has further balanced update age and data value under unstable connectivity \cite{cui2023data}, incorporated freshness into fairness-oriented scheduling \cite{zhu2024fairness}, required devices to be served within a bounded number of slots \cite{zhai2024uav}, and increased the priority of previously failed or unscheduled devices \cite{mestoukirdi2022uav}. These studies establish age-aware ranking and mechanisms to reduce starvation. FeRRy instead focuses on how these concerns interact with physical transport cost, where servicing an old or valuable update also requires flight time, contact time, sufficient radio reach, and remaining mission budget. 

Finally, asynchronous and hierarchical FL research addresses model-version mismatch when participant updates arrive at different times. Prior work has buffered asynchronous client updates before applying server-side aggregation, reducing sensitivity to individual stragglers while remaining compatible with secure aggregation \cite{nguyen2022federated}. Other hierarchical approaches perform asynchronous aggregation across both gateway and cloud tiers while jointly considering device selection and gateway association in heterogeneous IoT environments \cite{yu2023async}. These methods assume that updates are delivered through the network, whereas physical data muling introduces travel time as an additional source of staleness. FeRRy therefore treats aggregation as part of the mission lifecycle, where the same age and deadline information is used to determine whether a participant to be served also influences whether its transported update remains useful when the mule returns. The key distinction is thus the coupling of physical admission, service order, RF-dependent contact feasibility, and update freshness, rather than a new aggregation mechanism in isolation.

Table~\ref{tab:ferry-comparison} shows that prior work addresses individual parts of this problem, including FL-aware participant selection, freshness-aware scheduling, UAV trajectory planning, RF adaptation, and physical model transport. However, even the closest transporter-based approach does not address RF-dependent reach and service time with participant deadlines and in-mission route adaptation. FeRRy addresses this gap by using communication conditions, mobility, and FL-update urgency to plan and revise model transport under a finite mission budget.

\section{\sysname{} Methodology}
\label{sec:method}

% \sysname{} coordinates RF-aware UAV model transport under a finite mission budget. We describe the system workflow, the mission planner that selects the RF configuration and route, the in-flight adaptation of remaining contacts, and the aggregation of transported model updates.

% % -----------------------------------------------------------------------------
% \subsection{Setting, Layer Roles, and Decision Clocks}
% \label{sec:method:roles}

% \paragraph{Setting}
% A cluster has one edge server, one or more mules, and static field devices whose
% positions the server's registry holds. Each mule gets a disjoint slice
% $\mathcal{M}$ of the registry for the trial; with $K>1$ mules the slices are
% angular sectors around the shared dock, within one device of equal size, and the
% mules fly concurrently. Devices train offline between visits, so a contact is
% exchange only: the mule pushes the global model and pulls the device's prepared
% update. In \textsc{Pass 1} of a mission the mule collects updates along a route,
% merges them on board and docks to upload; the server merges across mules into
% the next global model. In \textsc{Pass 2} the mule delivers that model to its
% whole slice, which gives each later update a known basis version, and so an age.
% A budget $T_{\mathrm{miss}}$ bounds Pass~1: the plan and every check in flight
% require the predicted landing and upload within $T_{\mathrm{miss}}$ of takeoff.

% \paragraph{Learning task}
% The devices train a binary intrusion detector on CICIoT2023~\cite{CiCIoT2023}, a
% multilayer perceptron trained for one local epoch per model received, on 20{,}000
% balanced training flows split across them; a disjoint 6{,}000-flow test set
% scores each new global model, stamped when its closing upload completes. A device
% with no prepared update trains inside the contact, at no simulated cost. Every
% transfer is priced as $S=1$~MB.

% \paragraph{Layer roles and clocks}
% Table~\ref{tab:layer-roles} gives each layer's role. Hard gates (eligibility, the
% feasibility test, and the age cap) always run before any ranking, so a ranking
% rule, learned or not, chooses only among options the gates admit. On the
% \emph{plan clock}, once per mission at the dock, the mule chooses a contact band
% class $\bbar$ and a route together, because the band class fixes the contact
% range, the range fixes the stops, and the stops fix the route
% (Fig.~\ref{fig:clocks}). On the \emph{flight clock}, in Pass~1, it can revise the
% band and the next stop from the signal it observes, since interference varies
% over about one leg. Rate sets dwell, dwell advances mission time, mission time
% decides feasibility, and infeasibility triggers a re-plan. Learning can enter
% only through the flight slot.

% \begin{table}[t]
% \centering
% \caption{Role of each layer in the joint decisions.}
% \label{tab:layer-roles}
% \small
% \begin{tabularx}{\columnwidth}{l l L}
% \toprule
% \textbf{Layer} & \textbf{Role} & \textbf{Supplies to the decisions} \\
% \midrule
% L1 RF link & prices & Range, rate, dwell, and SNR over time for each band class; backhaul carrier \\
% L2 Scheduling & decides & Admission, band class $\bbar$, route, per-arrival (band, next stop), re-plan \\
% L3 Federated learning & values & Merge weight, update age, coverage weight, cross-mule merge \\
% \bottomrule
% \end{tabularx}
% \end{table}

% \begin{figure}[t]
% \centering
% \begin{tikzpicture}[
%   box/.style={draw, rounded corners=2pt, minimum height=0.62cm, align=center,
%               font=\scriptsize, inner sep=2.5pt},
%   gate/.style={box, fill=gray!15},
%   arr/.style={-{Latex[length=1.6mm]}, thick},
%   node distance=0.32cm]
%   % plan clock
%   \node[font=\scriptsize\bfseries, anchor=west] (pl) at (0,1.55) {Plan clock (dock, once per mission)};
%   \node[box] (d) at (0.9,0.95) {demand,\\weights};
%   \node[box, right=of d] (bb) {band class\\$\bbar$};
%   \node[box, right=of bb] (st) {stops at\\$R(\bbar)$};
%   \node[box, right=of st] (rt) {route\\$\pi$};
%   \node[gate, right=of rt] (g1) {feasibility\\test};
%   \draw[arr] (d)--(bb); \draw[arr] (bb)--(st); \draw[arr] (st)--(rt); \draw[arr] (rt)--(g1);
%   % flight clock
%   \node[font=\scriptsize\bfseries, anchor=west] (fl) at (0,0.1) {Flight clock (each Pass-1 arrival)};
%   \node[box] (ob) at (0.9,-0.5) {observe\\SNR per band};
%   \node[box, right=of ob] (pr) {(band, next\\stop) pair};
%   \node[gate, right=of pr] (g2) {feasibility\\test};
%   \node[box, right=of g2] (sv) {serve,\\dwell};
%   \node[box, right=of sv] (rp) {fits?\\else re-plan};
%   \draw[arr] (ob)--(pr); \draw[arr] (pr)--(g2); \draw[arr] (g2)--(sv); \draw[arr] (sv)--(rp);
%   \draw[arr, dashed] (rp.south) -- ++(0,-0.32) -| (ob.south);
%   \draw[arr, dashed] (g1.south) -- ++(0,-0.25) -| ($(ob.west)+(-0.15,0)$) -- (ob.west);
% \end{tikzpicture}
% \caption{The two decision clocks; the flight row shows FQ's per-stop choice.}
% \label{fig:clocks}
% \end{figure}

% % -----------------------------------------------------------------------------
% \subsection{Quantities Supplied by the Layers}
% \label{sec:method:models}

% \subsubsection{L1: Contact Link, Mission Time, and Backhaul}
% \label{sec:method:l1}

% \paragraph{Band classes}
% The contact link offers three band classes $\mathcal{B}$ on one carrier, LTE
% channel bandwidths~\cite{3gpp36104} with occupied bandwidth
% $B_b=N_{\mathrm{RB}}\times180$~kHz: \textsf{wide} (20~MHz, 100 resource blocks),
% \textsf{medium} (5~MHz, 25) and \textsf{narrow} (1.4~MHz, 6). With a common
% transmit power and a noise floor that scales with $B_b$, a narrower class
% reaches farther at a lower rate:
% \begin{equation}
% R_b^{3D} = R_0^{3D}\left(\frac{B_0}{B_b}\right)^{1/n},
% \qquad
% R_b = \sqrt{\left(R_b^{3D}\right)^2 - h^2},
% \label{eq:range}
% \end{equation}
% relative to a reference class $(B_0,R_0^{3D})$, with path-loss exponent $n=2.2$
% and altitude $h=25$~m. Anchoring \textsf{wide} at 60~m gives planar ranges of
% 60, 120 and 232~m, at which the mean SNR clears the decoding floor by a 5.1~dB
% shadowing margin $M_{\mathrm{sh}}$ (reach probability 0.9 under shadowing).

% \paragraph{Rate and dwell}
% The mean SNR at distance $d$ is
% $\bar\gamma_b(d)=\gamma_{\mathrm{floor}}+M_{\mathrm{sh}}+10\,n\log_{10}(R_b^{3D}/d^{3D})$,
% with $\gamma_{\mathrm{floor}}=-6.7$~dB (CQI~1), and the rate is
% \begin{equation}
% r_b(\gamma)=\min\bigl\{\kappa_b B_b\,\mathrm{SE}(\gamma),\;B_b\log_2(1+10^{\gamma/10})\bigr\},
% \label{eq:rate}
% \end{equation}
% for $\gamma\ge\gamma_{\mathrm{floor}}$ and zero below, with $\mathrm{SE}$ the CQI
% spectral efficiency and $\kappa_b$ an implementation efficiency. Moving $S_j$
% bytes takes $\tau_{j,b}^{\mathrm{dwell}}(\gamma)=8S_j/r_b(\gamma)$: $2S$ for a
% Pass-1 session, $S$ for a delivery or the upload, with the dwells at a stop
% adding up. The channel only decides whether a device is solicited (not if it is
% below the floor on arrival) and how long its session takes; a started session
% always completes, and the TCP messages carry no radio effects.

% \paragraph{Channel over time}
% The realized SNR of device $j$ on class $b$ at mission time $t$ is
% \begin{equation}
% \gamma_{b,j}(t)=\bar\gamma_b(d_j)+X_j(t)+A\sin\bigl(2\pi(t/P_c+\phi_b)\bigr)+\sigma_I\,\xi_b(t),
% \label{eq:realsnr}
% \end{equation}
% with correlated shadowing $X_j$ ($\sigma_{\mathrm{sh}}=4$~dB, 7.4~s correlation
% time) and per-class interference~\cite{3gpp36777} of period $P_c=60$~s, phases
% $\phi_b\in\{0,\tfrac13,\tfrac23\}$ and noise $\xi_b$; the interference-limited
% (``jittery'') regime sets $A=5$~dB and $\sigma_I=1.5$~dB. Keyed hashes give every
% policy the same draws. The planner prices only the mean $\bar\gamma_b$, so
% realized missions run longer than planned, and uses the outage probability
% \begin{equation}
% p_{\mathrm{out}}(b,d)=\Phi\!\left(\frac{\gamma_{\mathrm{floor}}-\bar\gamma_b(d)}{\sigma_{\mathrm{eff}}}\right),
% \quad
% \sigma_{\mathrm{eff}}^2=\sigma_{\mathrm{sh}}^2+\sigma_I^2+\tfrac{A^2}{2};
% \label{eq:pout}
% \end{equation}
% the flight clock observes $\gamma_{b,j}(t)$.

% \paragraph{Mission time and energy}
% Mission time is simulated time since takeoff, charged for transit at 5~m/s,
% dwell, a 1~s listen at a stop with a missing reply, the return, the upload and a
% 30~s dock turnaround. Propulsion energy follows the same ledger,
% \begin{equation}
% E = P_{\mathrm{move}}\,(t_{\mathrm{transit}}+t_{\mathrm{return}})
%   + P_{\mathrm{hover}}\,(t_{\mathrm{dwell}}+t_{\mathrm{listen}}),
% \label{eq:energy}
% \end{equation}
% with $P_{\mathrm{move}}=143.6$~W and $P_{\mathrm{hover}}=168.5$~W from Zeng et
% al.'s rotary-wing model~\cite{zeng2019energy}; no battery limit is enforced.

% \paragraph{Backhaul}
% The upload rides the best of three backhaul carriers and is lost with probability
% 2\%, without retry: Pass~2 then delivers the unchanged model, and the planner,
% which records the devices as merged before the upload, counts them as served.
% The component ablations' backhaul test and the radio study
% (Table~\ref{tab:studies}) use a time-varying backhaul instead, which loses about
% a fifth of the uploads on the fixed carrier (F closes 79\% of its rounds,
% against 97\%); only there does the adaptive carrier controller (F+L1) act, and
% since it sees the SNR that decides each loss, it is close to an oracle. The
% backhaul never enters a scheduling decision.

% \subsubsection{L2/L3: Demand, Deadlines, Ages, and the Age Cap}
% \label{sec:method:demand}

% \paragraph{Contact outcomes}
% A Pass-1 contact is \emph{clean} when the device returns an update with the
% expected round, size and checksum, \emph{partial} when it answers but refuses
% or fails those checks, and otherwise a \emph{timeout}. Whether a reachable device
% replies is drawn each mission from its reliability $\rho_j\sim U(0.15,1)$; a lost
% reply still costs the push. A session that outlasts its wall-clock timeout (36~s
% at $N=6$) is a timeout; results were identical at 0.75 to 1.5 times that value.
% A device can also be given a log-normal training time, restarted by each model
% it receives; an update not ready on arrival is a timeout, and the plan does not
% price training time.

% \paragraph{Per-device deadline}
% Each device $j$ has a fulfilment window $\Phi(j)$ and a deadline
% \begin{equation}
% \mathrm{Deadline}(j)=t+\Phi(j)-\iota(j),
% \label{eq:deadline}
% \end{equation}
% with $t$ the planning time and $\iota(j)$ the time since its last clean contact.
% After each contact the window is scaled and clamped,
% \begin{equation}
% \Phi(j)\leftarrow\min\{\Phi_{\max},\max\{\Phi_{\min},\,\beta_{o}\,\Phi(j)\}\},
% \label{eq:phi}
% \end{equation}
% with $\beta_{\mathrm{clean}}=0.8$; $\beta_{\mathrm{partial}}=1.25$ for a partial
% or a timeout after the device answered; and $\beta_{\mathrm{timeout}}=1.5$ for a
% silent, unreachable or left-out device. In units of $\sigma=\Tnom/10$~s, where
% $\Tnom$ is the median nominal two-pass mission on the wide band (203~s at $N=6$),
% $\Phi_0=60\sigma$, $\Phi_{\min}=5\sigma$ and $\Phi_{\max}=300\sigma$~s. A clean
% contact re-anchors $\iota(j)$; a failed one widens the window but leaves the
% anchor stale, so a device that keeps failing grows more overdue. A device the
% plan leaves out, or the mule drops in flight, counts as a timeout.

% \paragraph{Ages and coverage weight}
% The \emph{plan age} $a_j=m-U_j$ counts the missions since the mission $U_j$ whose
% merge last used the device's update, and strengthens its claim to be served. The
% \emph{merge age} $a_i=v-b_i$ of an update is the mule's model version minus the
% version $b_i$ the device trained from, which it echoes; it discounts the update
% in the merge. A miss relaxes a device's \emph{window} but not its
% \emph{priority}, which enters through the coverage weight
% \begin{equation}
% w_j=\max(a_j,1)\,(1+\mu_j),
% \label{eq:covweight}
% \end{equation}
% with $\mu_j$ the miss streak; under a plan, $\mu_j=a_j-1$, so $w_j\approx a_j^2$.

% \paragraph{Age cap}
% A device is \emph{capped} when $a_j\ge S$, with $S=2$ from a pilot (the fewest
% missions that serve every servable device on 90\% of layouts at the stress
% budget). A stop's deadline is the earliest among its uncapped members, and a
% stop of capped members only is exempt from it but not from the budget; a trimmed
% route flies stops with a capped member first; a capped device its stop cannot
% serve within the budget gets a one-device \emph{hover stop}; and the plan key
% first compares the ages of the capped devices a plan leaves out.

% % -----------------------------------------------------------------------------
% \subsection{The Feasibility Test}
% \label{sec:method:gate}

% One predicate decides whether a stop may be served next, and it is the only
% admission rule: it checks the committed plan, the plan search, every departure,
% the re-plan and the in-flight choices. For a stop $s$ with members
% $\mathcal{M}_s$, reached from a state (pose, mission time $c$, energy) on band $b$,
% \begin{align}
% t_{\mathrm{arr}} &= c+\|\mathrm{pose}-s\|/v,\\
% t_{\mathrm{fin}} &= t_{\mathrm{arr}}+\!\!\sum_{j\in\mathcal{M}_s}\!\tau^{\mathrm{dwell}}_{j,b}(\bar\gamma_b(d_j)),\\
% t_{\mathrm{home}} &= t_{\mathrm{fin}}+t_{\mathrm{return}}+t_{\mathrm{upload}},
% \end{align}
% and the stop is admitted only if $t_{\mathrm{fin}}$ meets its deadline,
% $t_{\mathrm{home}}$ falls within takeoff plus $T_{\mathrm{miss}}$, and the energy
% of Eq.~\eqref{eq:energy} fits any capacity set. The clauses are monotone in the
% member set, so a stop that fails whole can be reduced to the members that fit.
% The departure check catches a stop that the realized channel lengthens, but
% nothing re-checks after the last Pass-1 stop, so a mission can overrun.

% % -----------------------------------------------------------------------------
% \subsection{Plan Clock: Reach and Route}
% \label{sec:method:plan}

% At the dock the mule commits a plan $\mathcal{P}=(\bbar,\pi)$: the band class for
% both passes and an ordered list of stops, each reduced to the members that fit.

% \paragraph{Reach is a decision}
% For each class, a greedy anchor sweep clusters the slice into stops at the range
% $R_b$. A wide class serves few devices per stop at a high rate; a narrow class
% can reach the field from one stop but dwells longer. No fixed rule chooses
% between them, so the planner prices every class and commits the best.

% \paragraph{Route search}
% Routes pass the feasibility test stop by stop. Up to six devices, an exact search
% enumerates every ordered sequence of stops and member subsets and is optimal over
% what the score prices; with up to six stops, it searches ordered stop subsets;
% beyond that, a deterministic local search (2-OPT, then first-improvement moves,
% at most 2000 evaluations per class) runs. The empty plan is chosen only if no
% other is admitted.

% \paragraph{Plan score and plan key}
% Candidates are scored by
% \begin{equation}
% V(\bbar,\pi)=-\Bigl[c_1\Bigl(\tfrac{\Delta}{\Tnom}\Bigr)^{2}+c_2U+c_3L+c_4\frac{E}{P_{\mathrm{hover}}\Tnom}\Bigr],
% \label{eq:planscore}
% \end{equation}
% with $\Delta$ the duration of the two-pass mission on $\bbar$, $E$ its energy, and
% \begin{equation}
% U=1-\frac{\sum_{\text{served}}w_j}{\sum_{\text{demand}}w_j},
% \qquad
% L=\frac{\sum_{\text{served}}w_j\,p_{\mathrm{out}}(\bbar,d_j)}{\sum_{\text{demand}}w_j}
% \end{equation}
% the weighted coverage shortfall and expected link loss ($c_1=1$,
% $c_2=c_3=N_{\mathrm{demand}}$, $c_4=0.1$). The form follows FedEx's
% route-dependent bound~\cite{bian2025indirect}, but the constants are hand-set and
% the term in $\Delta$ is a staleness surrogate. Since every non-empty plan pays
% for a whole Pass~2, $V$ alone can prefer to fly empty, so the mule commits the
% plan with the smallest key
% \begin{equation}
% \bigl(\kappa_{\mathrm{cap}},\;-\textstyle\sum_{\text{served}}w_j/\sum_{\text{demand}}w_j,\;-V,\;\text{class index}\bigr),
% \label{eq:plankey}
% \end{equation}
% compared lexicographically, with $\kappa_{\mathrm{cap}}$ the ages of the capped
% devices the plan leaves out, largest first: time and energy break ties among
% plans that serve the same weight. A guard fold re-checks the route before
% commit, and every device left out is recorded as a miss.

% % -----------------------------------------------------------------------------
% \subsection{Flight Clock: Band, Next Stop, and Re-Planning}
% \label{sec:method:flight}

% \paragraph{The flight slot}
% In Pass~1 the mule can revise, per stop, the band on which to serve the stop it
% has reached and the next stop of the queue. F, the reference arm, leaves the
% slot empty and flies the plan's band and order. The fixed filling, FX, picks on
% arrival the fastest class that still reaches every device $\bbar$ reaches, at the
% observed SNR, and at the next departure, after the departure check, the nearest
% remaining stop whose move to the front passes the feasibility test from the
% realized mission time. The learned filling, FQ, makes both choices on arrival,
% as a pair $(b,s)$ ranked by a masked double DQN; a pair is admitted only if the
% band reaches what $\bbar$ reaches, the stop belongs to the plan, and the rest of
% the flight passes the feasibility test at the observed dwell, and with no
% admitted pair the mule serves on FX's band in the plan's order. Its features are
% the current and next-stop SNR per class, the dwell, deadline slack, plan ages,
% the mission time and energy left and the interference phase; it is trained in
% FerrySim on the merge weight collected less a time cost, with the shortfall of
% committed devices charged at the end of each sortie. FQ thus picks its next stop
% before serving, from the priced dwell, and FX after, from the realized time.

% \paragraph{Departure check and re-plan}
% After each stop, if the remaining queue no longer passes the feasibility test,
% the mule re-plans it: stops with a capped member move to the front, the rest
% keep the plan's order, and what no longer fits is trimmed and counted as a
% timeout for the mission.

% % -----------------------------------------------------------------------------
% \subsection{Aggregation: Mule, Cluster, and Delivery}
% \label{sec:method:merge}

% \paragraph{Merge on the mule}
% Devices send $\Delta\theta_i=\theta_i^{\mathrm{after}}-\theta_i^{\mathrm{basis}}$
% with its basis version, and at the end of Pass~1 the mule merges
% \begin{equation}
% w_i=n_i\,v_i\,s(a_i)\ \text{ for } a_i\le a_{\max,j},\qquad w_i=0\ \text{ otherwise},
% \label{eq:mergeweight}
% \end{equation}
% \begin{equation}
% M=\sum_{i\,\text{admitted}}n_iv_i,\qquad
% \Delta_m=\frac{1}{M}\sum_{i\,\text{admitted}}w_i\,\Delta\theta_i,
% \label{eq:mulemerge}
% \end{equation}
% with $n_i$ the local examples, $v_i$ a value proxy (uniform here), FedAsync's
% hinge $s(a)=1/(a+1)$~\cite{xie2019asynchronous}, and the cutoff
% $a_{\max,j}=\lfloor\Phi(j)/\Tnom\rfloor$, so the window that admits a device
% (Eq.~\eqref{eq:deadline}) also bounds how long its update counts. Dividing by $M$
% rather than $\sum w_i$ lets a uniformly stale mission move the model by only
% $s(a)$ times its mean; with current bases, as for almost every merge on the mule
% in our runs, it is the plain example-weighted mean.

% \paragraph{Merge across mules and delivery}
% At the server, each partial carries its own staleness $s_m=s(V-v_m)$, with $V$ the
% server's version and $v_m$ the mule's:
% \begin{equation}
% \theta\leftarrow\theta+\eta\,\frac{\sum_{m\,\text{live}}M_m\,s_m\,\Delta_m}{\sum_{m\,\text{live}}M_m},
% \label{eq:clustermerge}
% \end{equation}
% with $\eta=1$ and a partial live if non-empty. The server merges as soon as one
% partial arrives, so each upload closes its own round. A mule flies Pass~1 with
% the model of its previous dock, so with three mules at $N=6$, 160 of 198 merges
% applied a partial one to three versions old. If Pass~1 collected any update, the
% mule then delivers the new model to its slice, nearest stop first, on $\bbar$,
% skipping only a device below the floor.

% % -----------------------------------------------------------------------------
% \subsection{Mission Round, Cost, and Implementation}
% \label{sec:method:alg}

% \begin{algorithm}[t]
% \caption{One \sysname{} mission round}
% \label{alg:mission}
% \begin{algorithmic}[1]
% \Require Slice $\mathcal{M}$, model $\theta$ with version $v$, budget $T_{\mathrm{miss}}$, cap $S$
% \Statex \textbf{[Dock: plan clock]}
% \State Take slice $\mathcal{M}$ and the model $\theta$ (version $v$) received at the previous dock; compute demand, weights $w_j$, capped set
% \For{each band class $b\in\mathcal{B}$}
%   \State $\textit{stops}_b\gets$ cluster demand at $R_b$ and add hover stops for capped devices
%   \State $\mathcal{P}_b\gets$ \textsc{Search}$(\textit{stops}_b)$ under the feasibility test
% \EndFor
% \State $(\bbar,\pi)\gets$ plan with the smallest key (Eq.~\ref{eq:plankey}); guard fold; commit; record left-out devices as misses
% \Statex \textbf{[Pass 1: flight clock]}
% \For{each stop $s_k$ in the queue}
%   \State Fly to $s_k$; observe $\gamma_{b,j}(t)$ for all classes
%   \State Band $b_k$: $\bbar$ (F); fastest covering class (FX); with $s_{k+1}$, best admitted pair (FQ)
%   \State Serve members on $b_k$; update $\Phi(j)$ and $\iota(j)$ from each outcome
%   \If{remainder does not pass the feasibility test}
%     \State Re-plan by trimming; record dropped stops as timeouts
%   \EndIf
%   \State Next stop: plan order (F); nearest that keeps the remainder feasible (FX); the pair's (FQ)
% \EndFor
% \State Merge on the mule with cutoff (Eq.~\ref{eq:mulemerge}); fly to the dock; upload
% \Statex \textbf{[Dock: server]}
% \State Server merges across mules (Eq.~\ref{eq:clustermerge}) to obtain $\theta'$ with version $v+1$
% \Statex \textbf{[Pass 2, if Pass 1 collected any update]}
% \For{each stop in nearest-first order from the dock}
%   \State Deliver $\theta'$ on $\bbar$; devices resume offline training
% \EndFor
% \end{algorithmic}
% \end{algorithm}

% \paragraph{Cost and implementation}
% Planning cost is bounded (at most 1957 stop sequences per class for the exact
% search, 2000 evaluations for the local one), and each arrival evaluates at most
% $|\mathcal{B}|\times|\text{queue}|$ pairs; both merges are linear. The prototype
% runs the server, each mule and each device as separate Python processes over TCP,
% and the devices train the real model; mission time, the channel and energy are
% simulated. Updates carry integrity checks but no authentication;
% poisoning-robust aggregation and privacy are out of scope. The learned score, E3
% and the largest problems ($N=48$, $96$) run in FerrySim, an in-process simulator
% built from the same code that omits model training, flies one mule, uses an
% additive deadline law and no merge cutoff, and grows the field with $N$ in its
% scale cells.

FeRRy plans how a UAV collects and carries FL model updates when devices cannot
directly reach the edge server. We first describe the mission workflow, followed
by the RF and route planning, the decisions made during flight, and the model
aggregation process.

\subsection{System Model and Mission Workflow}
\label{sec:system-model}


We consider a federated learning deployment with one edge server, a set of
static field devices $\mathcal{D}$, and one or more UAV data mules. The server
maintains the current global model together with the location and participation
state of each registered device. Each device trains on its local data between
mule visits and stores its update until the next contact.

When multiple mules are used, we divide the devices into separate service
regions around a shared dock. We sort the devices by their angle from the dock,
starting after the widest empty arc, and divide them into contiguous sectors.
The sectors differ by at most one device, and each mule serves one sector for
the full trial. The mules operate at the same time and plan their routes
independently, while the edge server combines the updates that they return.

Each mule completes a two-pass mission. In \textsc{Pass~1}, the mule visits
selected contact locations, exchanges the current global model with reachable
devices, and collects their prepared local updates. The collected updates are
combined on the mule before it returns to the edge server. The server then
combines updates returned by different mules into the next global model. In
\textsc{Pass~2}, the mule delivers that model back to the devices in its region.
Each device stores the received model and its version so that the age of its next
update can later be determined.

We bound \textsc{Pass~1} by a mission budget $T_{\mathrm{miss}}$. A collection
plan must therefore include the time needed to travel between contacts,
exchange model data, return to the dock, and upload the collected updates. We
track these activities on the same mission clock so that device deadlines and
route feasibility use the same elapsed time.

The communication choice directly affects the route. A communication setting
with greater range may reduce the number of stops but increase transfer time.
A higher-rate setting shortens each contact but may require the mule to fly
closer to more devices. FeRRy therefore selects the communication setting as
part of the route-planning decision.

\subsection{RF-Aware Mission Planning}
\label{sec:mission-planning}

Let $\mathcal{B}$ denote the available RF classes. Each class $b$ has occupied
bandwidth $B_b$ and contact range $R_b$. For a reference class with bandwidth
$B_0$ and slant range $R_0^{3D}$, we model

\begin{equation}
R_b^{3D}
=
R_0^{3D}
\left(\frac{B_0}{B_b}\right)^{1/n},
\qquad
R_b =
\sqrt{\left(R_b^{3D}\right)^2-h^2},
\label{eq:contact-range}
\end{equation}

where $n$ is the path-loss exponent and $h$ is the UAV altitude. All RF classes
use the same transmit power, while the noise floor changes with occupied
bandwidth. Relative to the 20~MHz class, the 5~MHz and 1.4~MHz classes provide
about 6.0~dB and 12.2~dB more SNR, respectively. They can therefore satisfy the
same decoding threshold at a longer distance. The narrower classes provide
greater contact range, while the wider classes provide higher transfer rates.

For device $j$ and RF class $b$, the planner uses the expected mean SNR
$\bar{\gamma}_{b,j}$ to estimate the communication rate
$r_b(\bar{\gamma}_{b,j})$. For a payload of $S_j$ bytes, the expected
communication time is

\begin{equation}
\tau_{j,b}^{\mathrm{dwell}}
=
\frac{8S_j}{r_b(\bar{\gamma}_{b,j})}.
\label{eq:dwell}
\end{equation}

For each RF class, we use $R_b$ to group devices that can be served from the
same contact region and then construct candidate UAV stops. This produces a
different route for each communication choice. A class with greater range may
serve several devices from one stop, while a shorter-range class may require
separate stops. We represent a candidate mission as

\begin{equation}
\mathcal{P}=(b,\pi),
\end{equation}

where $b$ is the selected RF class and $\pi$ is the ordered sequence of contact
stops.

The planner also tracks how urgently each device should be served. Device $j$
maintains a service window $\Phi_j$ and a deadline

\begin{equation}
D_j=t+\Phi_j-\iota_j,
\label{eq:deadline}
\end{equation}

where $t$ is the planning time and $\iota_j$ is the elapsed time since its last
successful contact. After each contact, we update the service window as

\[
\Phi_j
\leftarrow
\min\{\Phi_{\max},
\max\{\Phi_{\min},\beta\Phi_j\}\}.
\]

We use $\beta=0.8$ after a successful contact, $1.25$ after a partial contact or
a timeout after the device responds, and $1.5$ when the device is silent,
unreachable, or not served. We initialize $\Phi_j$ with
$\Phi_0=60\sigma$ and limit it to
$\Phi_{\min}=5\sigma$ and $\Phi_{\max}=300\sigma$, where
$\sigma=T_{\mathrm{nom}}/10$. Here, $T_{\mathrm{nom}}$ is the median nominal
two-pass mission using the wide RF class.

\textcolor{red}{What exact age limit is used in the experiments, and what does the planner do differently once a device reaches it? Please give the threshold and the exact rule so we can describe it in one sentence.}
Service priority is tracked separately using

\begin{equation}
w_j=\max(a_j,1)(1+\mu_j),
\label{eq:coverage-weight}
\end{equation}

where $a_j$ is the number of mule missions since device $j$ last contributed an
update and $\mu_j$ is its consecutive miss count. A device therefore gains
priority when it has not contributed for several missions or has been missed
repeatedly. \textcolor{blue}{A device reaches the age limit when $a_j\ge S$, with $S=2$ missions in all
experiments; the planner then first minimizes the ages of such devices left
unserved, exempts stops containing only such devices from the deadline check but
not the budget check, keeps stops containing such a device first when trimming a
route, and gives a device that its stop cannot serve within the budget a separate
stop at its best hover point.}

To determine whether a stop can be added to the mission, suppose the mule is at
position $p$ at mission time $t$ and considers stop $s$ with device set
$\mathcal{D}_s$. We estimate

\begin{align}
t_{\mathrm{arr}}
&=
t+\frac{\|p-s\|}{v},
\\
t_{\mathrm{fin}}
&=
t_{\mathrm{arr}}
+
\sum_{j\in\mathcal{D}_s}
\tau_{j,b}^{\mathrm{dwell}},
\\
t_{\mathrm{home}}
&=
t_{\mathrm{fin}}
+
\frac{\|s-p_{\mathrm{dock}}\|}{v}
+
\tau_{\mathrm{upload}},
\end{align}

where $v$ is the UAV cruise speed. We add the stop only if its device deadlines
can still be met and the mule can return to the dock and upload the collected
updates before the \textsc{Pass~1} budget expires.

\textcolor{red}{How exactly is the expected communication-loss cost \(L\) calculated? We define what it represents, but not how it is computed. Please provide the formula or implementation rule.}
FeRRy first favors plans that serve devices with the highest service need.
\textcolor{blue}{Among plans that serve the same weight,} it uses

\begin{equation}
V(b,\pi)
=
-
\left[
c_1\left(\frac{\Delta}{T_{\mathrm{nom}}}\right)^2
+
c_2 U
+
c_3 L
+
c_4\frac{E}{P_{\mathrm{hover}}T_{\mathrm{nom}}}
\right],
\label{eq:mission-score}
\end{equation}

where $\Delta$ is the expected mission duration, $U$ is the weighted demand left
unserved, $L$ is the expected communication-loss cost, and $E$ is the expected
flight and hovering energy. We use this expression only to compare candidate
plans.

{\color{blue}
The coverage and loss terms are
\begin{equation}
U=1-\frac{\sum_{j\in\mathcal{S}}w_j}{\sum_{j\in\mathcal{D}_{\mathrm{dem}}}w_j},
\qquad
L=\frac{\sum_{j\in\mathcal{S}}w_j\,p_{\mathrm{out}}(b,d_j)}
       {\sum_{j\in\mathcal{D}_{\mathrm{dem}}}w_j},
\label{eq:coverage-loss}
\end{equation}
where $\mathcal{S}$ is the set of served devices, $\mathcal{D}_{\mathrm{dem}}$ the
devices requesting service, and $d_j$ the planar distance from device $j$ to its
stop. The outage probability
\begin{equation}
p_{\mathrm{out}}(b,d)=
Q\!\left(\frac{\bar{\gamma}_b(d)-\gamma_{\min}}{\sigma_{\mathrm{eff}}}\right),
\qquad
\sigma_{\mathrm{eff}}^2=\sigma_s^2+\sigma_I^2+\frac{A^2}{2},
\label{eq:outage}
\end{equation}
treats the SNR as Gaussian around its mean, where $Q$ is the Gaussian tail
function and $\gamma_{\min}=-6.7$~dB is the decoding floor. The parameters
$\sigma_s$, $\sigma_I$, and $A$ are defined in Section~\ref{sec:environment}; at
their defaults, $\sigma_{\mathrm{eff}}\approx5.5$~dB.\par}

We set $c_1=1$, $c_2=c_3=N_{\mathrm{demand}}$, where
$N_{\mathrm{demand}}$ is the number of devices requesting service, and
$c_4=0.1$. These values are fixed rather than fitted from the evaluation data.
The larger values of $c_2$ and $c_3$ give uncovered demand and communication
loss greater weight, while the energy term has a smaller weight because it
largely follows mission duration.

When at most six devices request service, we enumerate all ordered stop
sequences and device groups. For larger cases with at most six candidate stops,
we search ordered subsets of those stops. For larger search spaces, we use
2-OPT followed by first-improvement moves, with at most 50 passes or 2,000
candidate evaluations for each RF class. The selected plan determines the RF
class and initial contact order before takeoff.

\subsection{In-Mission Adaptation}
\label{sec:inmission}

The mission plan is based on expected channel conditions, but the channel
observed during flight may differ. FeRRy therefore checks the remaining mission
after each \textsc{Pass~1} stop using the mule's current position, elapsed time,
and measured SNR.

At stop $s_k$, the mule selects the fastest RF class that can still reach the
devices assigned to that stop. This class may differ from the one selected
before takeoff. After the contact finishes, the observed communication time
updates the mission clock, and FeRRy applies the feasibility test from
Section~\ref{sec:mission-planning} to the remaining route.

FeRRy also checks whether another remaining stop should be visited next. It
selects the nearest stop that can be moved forward without making the rest of
the mission infeasible. If no such change improves the route, the mule keeps the
original order. We represent this decision as

\begin{equation}
a_t=(b_t,s_{t+1}),
\end{equation}

where $b_t$ is the RF class used at the current stop and $s_{t+1}$ is the next
contact location.

A contact that takes longer than expected may also leave insufficient time for
the original route. After each stop, FeRRy therefore checks the full remaining
sequence. If the route no longer fits, it first keeps stops containing devices
with the highest service need and removes contacts that cannot be completed
within the remaining mission time. Devices that are removed are recorded as
missed and receive greater priority in later missions.

We also evaluate a learned alternative that ranks feasible
$(b_t,s_{t+1})$ choices using the current RF and mission state. The same
reachability, deadline, and mission-budget checks are applied before this
ranking. The learned policy therefore changes only the order among choices that
are already feasible. Section~\ref{sec:evaluation} compares this approach with
the deterministic rule and with a plan that is followed without adaptation.

\subsection{Model Aggregation and Delivery}
\label{sec:aggregation}

At the end of \textsc{Pass~1}, each mule combines the updates collected during
its mission before returning to the edge server. Device $i$ sends its local
update $\Delta\theta_i$ together with the version of the global model used for
training. This version is used to determine the age of the update and reduce
the contribution of older updates when necessary.

\textcolor{red}{Please provide the actual staleness function \(s(a_i)\) and the age cutoff used to exclude an update. Are the same function and cutoff used on both the mule and server?}

For the set of accepted updates $\mathcal{A}$, the mule computes

\begin{equation}
\Delta_m =
\frac{1}{\sum_{i\in\mathcal{A}} n_i}
\sum_{i\in\mathcal{A}}
n_i\,s(a_i)\,\Delta\theta_i ,
\label{eq:mule-aggregation}
\end{equation}

where $n_i$ is the number of local training samples, $a_i$ is the update age,
and $s(a_i)$ gives less weight to older updates. Updates that exceed the allowed
age are not included. When all collected updates are current, this becomes
standard example-weighted averaging. \textcolor{blue}{We use FedAsync's staleness function $s(a)=1/(a+1)$~\cite{xie2019asynchronous},
where $a_i$ is the mule's model version minus the version device $i$ trained from.
The mule excludes an update when $a_i>\lfloor\Phi_i/T_{\mathrm{nom}}\rfloor$, so
the service window that admits a device also limits how long its update counts.}

When several mules return at different times, the server applies the same
age-based weighting to their updates. A later arrival may have been trained
from an older global model, so its contribution is reduced according to the
difference in model version. Valid updates are incorporated as they arrive
rather than waiting for every mule. \textcolor{blue}{The server uses the same function on $V-v_m$, where $V$ is its model version and
$v_m$ the version the mule carried, but applies no age cutoff.}

After aggregation, the new global model is distributed during
\textsc{Pass~2}. Each device stores the model and its version before continuing
local training. That version is then used to determine the age of the device's
next update.

\section{Evaluation Methodology}
\label{sec:evaluation}
We evaluate \sysname{} using a prototype with model training and a simulator
for learned decisions and larger deployments. Both use the same mission planner, feasibility test, RF model, and mission logic.

\subsection{Prototype and FL Workload}
% Real prototype vs. FerrySim, CICIoT2023, model, payload.
\label{sec:prototype}

\textcolor{red}{Please give the exact FerrySim differences needed for reproducibility, what additive deadline rule does it use, and how does the field size grow with \(N\) in the scalability experiments?}

Our prototype runs the edge server, each UAV mule, and each field device as
separate Python processes that communicate over TCP. Devices train and exchange
real model parameters. UAV movement, wireless conditions, mission time, and
propulsion energy are simulated so that the same conditions can be repeated
across trials.

We use CICIoT2023~\cite{CiCIoT2023} as the FL workload and formulate binary
intrusion detection by grouping the attack classes into one attack label. Each
trial uses 20,000 training flows and a separate test set of 6,000 flows. The
training data are divided across the participating devices, and each device
performs one local epoch after receiving a new global model. Unless otherwise
stated, the main experiments use an IID partition; Section~\ref{sec:robustness}
also evaluates non-IID data.

The local model is a multilayer perceptron with hidden layers of 64, 32, 16, 8,
and 4 units, trained with binary cross-entropy and Adam. The trained model is
approximately 18.8~KB. For the main experiments, we use a 1~MB transfer size to
represent a larger FL model. At 18.8~KB, transfers within 60~m take less than
0.2~s on all RF classes, leaving little difference between their rates. At
1~MB, transfer time ranges from about 1 to 8~s, making both communication range
and transfer rate relevant. Section~\ref{sec:robustness} also evaluates the
measured 18.8~KB model and payloads of 2, 5, and 10~MB.

For learned decisions and larger populations, we use FerrySim, a simulator that
uses the same planning and mission logic as the prototype without running model
training or communication processes. This allows us to train the learned policy
and study populations of up to 96 devices. \textcolor{blue}{Devices do not train in FerrySim: each holds an equal data share, so it reports
served devices and returns rather than model accuracy. It also differs from the
prototype in four ways. It uses one mule. It applies an additive deadline update,
$\Phi_j\leftarrow\max\{\Phi_j-5\sigma,\,5\sigma\}$ after a successful contact and
$\Phi_j\leftarrow\Phi_j+10\sigma$ after any miss, with no upper limit. It applies
$s(a)$ but no merge cutoff. Its scale experiments use a field half-width of
$100\sqrt{N/6}$~m, giving 200, 283, and 400~m at $N=24$, 48, and 96 at the device
density of $N=6$, with knee budgets of 350, 850, and 1,660~s. Its $N=6$ and $N=12$
settings keep the prototype's field.} We identify FerrySim results separately.

\subsection{Network and Mission Environment}
\label{sec:environment}

We place devices uniformly at random in a \(200\times200\)~m area centered on
the UAV dock and generate a new layout for each trial. The prototype experiments
use \(N\in\{6,12,24\}\) devices. We use one mule as the main setting and up to
three mules to study parallel operation.

\textcolor{red}{Since the 60 m wide-band range is a simulation assumption chosen so that range affects the route, do we have any sensitivity experiment using another reference range or field size? If yes, please provide it. If not, please confirm so we can state this clearly as an evaluation limitation.}

We use three LTE channel bandwidths: 20, 5, and 1.4~MHz. With UAV altitude
\(h=25\)~m, path loss exponent \(n=2.2\), and a 60~m reference range for the
20~MHz class, the corresponding contact ranges are approximately 60, 120, and
232~m. The 60~m reference is a simulation setting chosen so that communication
range affects route planning within the \(200\times200\)~m field; it is not a
testbed measurement. Wider classes provide higher transfer rates over shorter
distances, while narrower classes provide greater range.

\textcolor{blue}{We did not vary the 60~m reference, which is a limitation of this evaluation. The
narrow class can reach the whole \(200\times200\)~m field from a single stop at
the dock. As exploratory checks at
\(N=6\) and the knee budget, with 20 paired trials, we varied the narrow reach and
device placement. Reducing the narrow reach from 232~m to 120 and 180~m raises
\sysname{}-Plan's time to target from 194~s to 214 and 209~s, still significantly
faster than every baseline (265--282~s). With 25\%, 50\%, and 75\% of devices
placed beyond 60~m, \sysname{}-Plan needs 150, 180, and 209~s, against 221--233,
190--228, and 245--281~s. All six baselines are significantly slower at 25\%, none
at 50\%, and three of six at 75\%.}

\textcolor{blue}{We also flew 12 devices on \(400\times400\) and \(600\times600\)~m fields, which
the narrow class cannot cover from the dock, with 20 paired trials per cell and a
900~s knee budget on each field. With a 1~MB model, \sysname{}-Plan needs 1,000 and
1,534~s, significantly faster than Deadline-Heuristic, MAX-AoI, Oort, and
Value-Aware FL (by 321--831~s) but not FedEx. With a 10~MB model, it selects the
wide class on every mission, as the baselines do, and needs 1,434 and 1,931~s, with
no significant difference from FedEx or FedCS, with or without the budget. With 24
devices on the \(600\times600\)~m field, a 10~MB model, and a 600~s budget, no
baseline differs significantly from \sysname{}-Plan, at 3,361~s against
2,938--3,550~s. No baseline is significantly faster in any of these cells, but the
lead over FedEx at \(N=6\) does not carry over to these fields.}

We vary the channel during each mission so that the conditions observed in
flight can differ from those expected during planning. For device \(j\) and RF
class \(b\), the realized SNR is

\begin{equation}
\gamma_{b,j}(t)
=
\bar{\gamma}_{b}(d_j)
+
X_j(t)
+
A\sin\!\left(2\pi\left(\frac{t}{P_c}+\phi_b\right)\right)
+
\sigma_I\xi_b(t),
\label{eq:channel-model}
\end{equation}

where \(\bar{\gamma}_{b}(d_j)\) is the mean SNR at distance \(d_j\),
\(X_j(t)\) represents correlated shadowing, and the remaining terms represent
time varying interference. In the default setting, we use 4~dB shadowing,
\(A=5\)~dB, \(\sigma_I=1.5\)~dB, and \(P_c=60\)~s. The planner uses the
mean SNR, while the decisions made during flight use the observed SNR. We vary
interference, path loss, and shadowing in the robustness experiments.

The UAV flies at a fixed altitude with a cruise speed of 5~m/s. Mission time
includes travel, communication time, a 1~s listening period when an expected
response is missing, return travel, model upload, and a 30~s dock turnaround.
We estimate propulsion energy using
\(P_{\mathrm{move}}=143.6\)~W during flight and
\(P_{\mathrm{hover}}=168.5\)~W during communication and listening. We record
energy as a mission cost but do not impose a battery limit.

For each device population, we select two \textsc{Pass~1} mission budgets from
pilot experiments. The \emph{knee} budget is the smallest budget whose mean
update yield reaches 95\% of the best yield observed in the pilot runs. The knee
budgets are 150, 180, and 262~s for \(N=6,12,\) and \(24\), respectively.
The \emph{stress} budget is half of the corresponding knee budget. For the
\(N=6\) scheduler comparison, we also evaluate a 300~s budget and exploratory
budgets of 60 and 45~s.

The mule returns its collected update over a separate backhaul link. In the main
setting, it selects the best of three backhaul carriers, and an upload is lost
with probability 2\%. A time varying backhaul is used only in the corresponding
component and channel experiments.

\subsection{Compared Approaches}
\label{sec:baselines}

\textcolor{red}{We need a short description of the learned-policy setup. Please provide the network architecture, state/features, action space, reward, training duration/episodes, optimizer/learning rate, and whether FeRRy-RL and Model-Aided RL received comparable training.}

We compare \sysname{} with scheduling, freshness, UAV routing, and physical
model transport approaches that represent the closest alternatives to parts of
our design. All methods use the same device layouts, channel conditions, mission
budgets, and paired random seeds. Table~\ref{tab:arms} summarizes the methods.

We evaluate three \sysname{} variants. \textbf{\sysname{}-Plan} follows the RF
class and route selected before takeoff. \textbf{\sysname{}-Adapt} adds the
deterministic decisions described in Section~\ref{sec:inmission}.
\textbf{\sysname{}-RL} uses a learned policy to rank the feasible RF class and
next stop choices. These variants separate the effect of the initial mission
plan from deterministic and learned decisions during flight.

\textcolor{blue}{\sysname{}-RL and Model-Aided RL are double DQNs trained in FerrySim with the same
network, two hidden layers of 64 tanh units giving one Q-value per candidate, and
the same training setup: Huber loss, batches of 64, a 50,000-transition replay
buffer, target updates every 500 steps, and at most 10,000 episodes with early
stopping. \sysname{}-RL scores each feasible $(b_t,s_{t+1})$ pair from 36 features
covering SNR, dwell, travel, deadline slack, age, remaining time, and recent
channel state. Its reward is the merge weight collected at a stop, less 0.1 times
the time spent over $T_{\mathrm{nom}}$, with a coverage penalty at the end of the
mission. It uses Adam at $10^{-3}$, with ten seeds at each of six discount factors
from 0 to 0.99 on \(N=6\) and \(N=12\) settings; we fly the median-validation seed
at the best factor, 0.75. It trains under FerrySim's additive deadline update but
flies under the update of Section~\ref{sec:mission-planning} in the prototype.
Model-Aided RL scores each remaining stop from ten features adapted from
Chen~\cite{chen2023model}. Its reward is the share of updates collected at the
stop, and it follows the published settings: Adam at $5\times10^{-4}$, discount
0.99, and exploration from $\epsilon=1$, over five seeds on the wide class. The two
therefore share the simulator, network, and training limit, but not the reward,
learning rate, or seed count.}

\begin{table}[t]
\centering
\caption{Approaches compared in the evaluation.}
\label{tab:arms}
\small
\begin{tabular}{p{0.25\columnwidth} p{0.68\columnwidth}}
\hline
\textbf{Approach} & \textbf{Description} \\
\hline

\textbf{\sysname{}-Plan} &
Selects the RF class and route before takeoff using device deadlines, coverage,
and the mission budget. \\

\textbf{\sysname{}-Adapt} &
Adds deterministic changes to the RF class and next contact during flight. \\

\textbf{\sysname{}-RL} &
Uses a learned policy to rank feasible RF class and next stop choices. \\

\textbf{Live-FL} &
No-mule reference in which FL uses the degraded infrastructure link. \\

\textbf{Deadline-Heuristic} &
Our earlier deadline based mule scheduler without joint RF and route planning. \\

\textbf{MAX-AoI} &
Serves participants according to age of information~\cite{kadota2018scheduling}. \\

\textbf{Oort} &
Selects participants using statistical utility and participation
history~\cite{lai2021oort}, adapted to the mule setting. \\

\textbf{Value-Aware FL} &
Ranks participants using data value, update age, and expected
connectivity~\cite{cui2023data}. \\

\textbf{FedEx} &
Uses physical model transport with a visit-all 2-OPT
route~\cite{bian2025indirect}. \\

\textbf{FedEx-Merge} &
Uses the FedEx route together with its original aggregation rule. \\

\textbf{FedCS} &
Uses deadline based participant selection following
FedCS~\cite{nishio2019client}, adapted to the disconnected setting. \\

\textbf{Model-Aided RL} &
Uses a learned next stop policy adapted from
Chen~\cite{chen2023model}, without FeRRy's RF and route plan. \\

\hline
\end{tabular}
\end{table}

\textcolor{red}{What exact device-assignment rule does FedEx use when there are multiple mules? We currently only say that it uses its own assignment rather than FeRRy's angular sectors.}
The comparison methods keep their main scheduling principle while operating in
the same mule setting. MAX-AoI, Oort, and Value-Aware FL rank candidate
contacts using their respective age or utility measures. FedCS selects
participants under a round deadline. Deadline-Heuristic uses device deadlines
and distance to construct its route. FedEx follows a visit-all route and does
not remove devices when the route exceeds the mission budget. \textcolor{blue}{With several mules, FedEx assigns devices using its CARP rule~\cite{bian2025indirect}:
annealed Gibbs sampling minimizes $\sum_k R_k T_k^2$, where $R_k$ is the number of
devices assigned to mule $k$ and $T_k$ is its round-trip time, the 2-OPT tour time
plus $R_k$ transfer times. We compute the assignment once per trial from the device
positions and omit the original energy constraint.} Model-Aided RL selects each
next contact from the stops that remain feasible.

In the main comparison, these methods use the wide RF class and therefore do
not choose communication range as part of their scheduling decision.
\sysname{}-Plan and its variants instead search across the available RF classes.
We separately evaluate fixed RF classes to isolate the effect of communication
range from the remaining planning decisions.

\subsection{Metrics and Statistical Analysis}
\label{sec:metrics}

Our primary metric is the simulated time required for the global model to reach
a target accuracy \(\tau\), measured from the first UAV takeoff to the first
server evaluation that reaches the target. We set \(\tau=0.71\) before the main
experiments. This is the largest value on a 0.01 grid reached by at least 80\%
of the Deadline-Heuristic trials at every knee budget. The target is above the
0.667 accuracy obtained by predicting every test sample as benign. We also
report final model accuracy when some methods do not consistently reach
\(\tau\).

We use several mission metrics to explain the time to target. \emph{Update
yield} is the number of device updates included in each round.
\emph{Round-close rate} is the fraction of missions that produce a new global
model. We measure update freshness as the number of mule missions since each
device last contributed an update. We also report missed deadlines, mission
budget overruns, and propulsion energy where relevant.

All methods within an experiment use paired random seeds and therefore see the
same device placement, channel conditions, and data partition. We report the
paired mean difference with a 95\% percentile bootstrap confidence interval
using 2,000 resamples. We test paired differences using the Wilcoxon signed-rank
test and apply Holm correction within each comparison group.

For time to \(\tau\), the paired comparison uses trials in which both methods
reach the target, and we report the target reach rate separately. We consider a
difference statistically supported when the confidence interval excludes zero
and the Holm corrected \(p\)-value is below 0.05. The comparison groups,
methods, and primary metrics were specified before each study. Exploratory
experiments are marked with \(^{\S}\) and are reported separately from these
tests. The main scheduler and component studies use 40 paired trials per
setting; the remaining studies use 20 unless stated otherwise.

% \section{Experimental Setup and Design}
% \label{sec:setup}

% \subsection{Testbed}
% \label{sec:setup:testbed}

% Each trial flies four missions of the full system with a 1~MB payload, in the
% interference-limited (jittery) regime, and a separate scorer computes the
% metrics from the events it records. FerrySim trains the learned per-stop score
% and the learned baseline E3, and runs the largest problems ($N=48$ and $96$).

% \subsection{Cells, Budgets, and the Accuracy Target}
% \label{sec:setup:cells}

% A cell fixes the number of devices $N$, the number of mules $K$, and the mission
% budget. Devices are static and placed uniformly at random on a $200\times200$~m
% square centred on the dock, a new layout per trial; the square is the same at
% every $N$, so the field grows denser with $N$. Pilots set two budgets per $N$
% before the studies ran: the \emph{knee}, the smallest budget whose mean update
% yield reaches 95\% of the best on a grid (150, 180 and 262\,s at $N=6$, 12 and
% 24), and the \emph{stress} budget, half of it. The accuracy target $\tau=0.71$
% is the largest value, on a 0.01 grid, that at least 80\% of H1's trials reach at
% every $N$'s knee. It lies 4.3 points above the 0.667 of labelling every test flow
% benign, so time to $\tau$ measures how soon a model becomes useful, not how good
% it ends.

% \subsection{Arms}
% \label{sec:setup:arms}

% Table~\ref{tab:exp5_arms} lists the arms: \sysname{}'s three fillings of the
% flight slot (F, FX, FQ) and the baselines, each a whole scheduler flown in the
% same simulator, at the same budgets, on the same seeds and on the wide contact
% band, with its departures from its paper stated.

% \begin{table}[t]
% \centering
% \caption{Arms compared in the evaluation.}
% \label{tab:exp5_arms}
% \small
% \begin{tabularx}{\columnwidth}{l X}
% \toprule
% \textbf{Arm} & \textbf{Description} \\
% \midrule
% F & FeRRy: dock-time plan (band class, route, deadline-gated admission, coverage term, age cap $S$), committed in flight. \\
% FX & F with the cross-layer in-flight rule: at each stop, the nearest stop that keeps the plan feasible, at the fastest band that still reaches the planned targets. \\
% FQ & F with the learned in-flight rule: FX's slot, masks and fallback, ranked by a learned pair score (masked double DQN trained in FerrySim, $\gamma=0.75$). \\
% H1 & Staged deadline heuristic (our earlier scheduler): deadline-tiered contact regions, no dock-time plan. \\
% H0 & Live-link reference: FL over the degraded infrastructure link, no mule. \\
% D1 & MAX-AoI: visit the stalest device next \cite{kadota2018scheduling}. \\
% D2 & Oort-style statistical utility \cite{lai2021oort}. \\
% D3 & Whittle index on value-weighted age of updates \cite{cui2023data}, our expected-connectivity variant. \\
% D4 & FedEx's visit-all 2-OPT tour \cite{bian2025indirect}, one transporter, never skips; D4$_{\text{FedEx}}$ adds its $1/N$ delta merge. \\
% D5 & FedCS, degraded: round-deadline selection without its resource-request phase \cite{nishio2019client}. \\
% E3 & Single-agent next-stop DQN after Chen et al.\ \cite{chen2023model}, trained in FerrySim. \\
% \bottomrule
% \end{tabularx}
% \end{table}
% \paragraph{How the baselines fly}
% The baselines share \sysname{}'s clustering at the wide band's range, its
% feasibility test and its unbudgeted, nearest-first Pass~2, but they serve whole
% stops only, have no age cap, and use an additive deadline law: a device's window
% shrinks by $5\sigma$~s after a clean contact and widens by $10\sigma$~s after a
% miss, with no ceiling. H1 admits stops in deadline order under the deadline and
% budget clauses, flies new devices first and then by distance from the dock, and
% trims at each departure. D1--D3 walk their own ranking (age of information,
% statistical utility, the Whittle index) under the budget alone, and D5 adds
% stops by least marginal time while they fit one round deadline; in flight, these
% four re-check only the budget and re-admit by their own rule. D4 flies FedEx's
% visit-all tour with no budget check, and with three mules splits the devices by
% FedEx's own assignment. E3 picks each next stop by its learned Q-values among the
% stops that still fit the budget, with no re-plan; it is trained with one mule
% (learning rate $5\times10^{-4}$, $\gamma=0.99$) and flown per mule with three.

% \subsection{Metrics and Statistics}
% \label{sec:setup:metrics}

% The primary metric is the simulated time to $\tau$, from the first takeoff to the
% first global model that reaches it. We also report the network age of updates
% (the mean number of missions since each device's last merged update), the
% round-close and deadline-miss rates, the update yield (updates merged per
% round), and mission energy. Trials are paired by seed, so every arm in a cell
% sees the same layout, channel and data. Each comparison is the paired mean
% difference with a 95\% percentile-bootstrap CI (2000 resamples) and a Wilcoxon
% signed-rank $p$, Holm-corrected within each study; time to $\tau$ is compared
% over the seeds in which both arms reach it. A comparison is a \emph{claim} when
% the CI excludes 0 and the Holm $p<0.05$.

% \subsection{Experiment Design}
% \label{sec:setup:design}

% The evaluation answers six questions: Q1, does \sysname{} beat the state of the
% art; Q2, where does its advantage come from; Q3, how does it scale, and what does
% deciding cost; Q4, does each part of its design matter; Q5, does learning help
% inside its plan; and Q6, how robust is it. Table~\ref{tab:studies} lists the
% studies that answer them. An analysis plan fixed each study's cells, arms,
% primary metric and test family before it ran; unless stated, a study flies one
% mule with 20 paired trials per cell. Q4 tests five design claims, each against an
% alternative that removes or replaces one part of the design: C1, reach is a
% decision; C2, one derived objective; C3, one deadline in three roles; C4, two
% clocks, re-decided per stop; and C5, fairness under physical cost. The
% learned-score study keeps the learned score only if some $\gamma$ beats
% $\gamma=0$ by 0.01 in held-out return (ten seeds per $\gamma$, 1000 episodes per
% cell); the FQ-vs-E3 comparison, added after that verdict, is exploratory.

% \begin{table*}[t]
% \centering
% \caption{The studies of the evaluation. $^{\dagger}$Flies the time-varying
% backhaul.}
% \label{tab:studies}
% \footnotesize
% \begin{tabularx}{\textwidth}{@{}r l c L l l@{}}
% \toprule
% \textbf{\#} & \textbf{Study} & \textbf{Q / C} & \textbf{Compares} & \textbf{Cells} & \textbf{Primary metric} \\
% \midrule
% 1 & Merge rules & Q4, C3 & Plain mean, FedBuff, Async-HFL and a FedProx term against \sysname{}'s merge & $N{=}6$; knee, stress & time to $\tau$ \\
% 2 & Deadline forms & Q4, C3 & FedCS's round deadline, Oort's preferred duration, an additive law & $N{=}6$; knee, stress & deadline-miss rate \\
% 3 & Scheduler comparison & Q1 & F and FX against H1, D1--D5 and E3; H0 described & $N{=}6$; $K{=}1,3$; knee, stress & time to $\tau$ \\
% 4 & Band pinning & Q2, C1 & F pinned to each band class; an exhaustive-plan oracle in FerrySim & $N{=}6$; knee; 1~MB, 18.8~KB & age of updates \\
% 5 & Learned score & Q5, C4 & FerrySim: $\gamma$ from 0 to 0.99 against FX and a greedy rule; full system: FQ against FX and F & FerrySim; $N{=}12$ & held-out return; time to $\tau$ \\
% 6 & Objective terms & Q4, C2 & FX without its coverage term or its dwell term; D4 & $N{=}12$; knee, stress & round-close rate \\
% 7 & Fairness & Q4, C5 & MAX-AoI, Whittle index, FedEx route; F without coverage term, age cap or miss priority & $N{=}6$; knee, stress; 4 missions & age of updates \\
% 8 & Scale & Q3 & F, FX, H1, D3, D4 as the field fills (FQ at $N{=}12$) & $N{=}6,12,24$; knee, stress & time to $\tau$ \\
% 9 & Decision cost & Q3 & (a) planner wall time; (b) one to three mules; (c) FerrySim budget sweeps & $N{=}6$--$96$; $K{\le}3$ & plan time; time to $\tau$ \\
% 10 & Training time & Q6 & Log-normal training time (median 68\,s), and 20\% of devices $5\times$ slower; F, FX, H1, D4, D5 & $N{=}6$; knee & time to $\tau$ \\
% 11 & Non-IID data & Q6 & Dirichlet split ($\alpha{=}0.1$) over attack families; F, FX, H1, D3 & $N{=}6$; knee & time to $\tau$ \\
% 12 & Component ablations & Q4 & F with one switch flipped (whole stops, local search, cap, hover, adaptive backhaul$^{\dagger}$) & $N{=}6$; knee, stress; 40 trials & age of updates \\
% 13 & Radio$^{\dagger}$ & Q6 & Interference 1 or 17\,dB, path-loss exponent 2.7, shadowing 8\,dB; F, FX, H1 & $N{=}6$; knee & time to $\tau$ \\
% \midrule
% -- & FQ vs E3 (exploratory) & Q5 & FQ against E3, F, FX and D4 & $N{=}6$; knee, stress & time to $\tau$ \\
% \bottomrule
% \end{tabularx}
% \end{table*}

% \paragraph{What was run}
% To meet the submission date, the merge-rule, deadline-form and objective-terms
% studies, the learned score's full-system check, and the scheduler comparison's
% D5 and E3 arms flew 20 of their planned 40 trials per cell; the remaining trials
% extend the same seeds. We dropped the merge-rule study's H1 route with an
% unbudgeted Pass~2, the scheduler comparison's relaxed budget, the band-pinning
% study's range and far-device settings, the fairness study's eight-mission
% horizon, the scale study's multi-mule and clean-channel cells, the training-time
% study's other payloads, the non-IID study's other partitions, and the radio
% study's adaptive-backhaul arms. Two planned studies, on when learning helps and
% on real radios (which awaits testbed access), did not run. FQ left the scheduler
% comparison and the scale study, as the plan specified, when the learned-score
% study did not keep the learned score.
\begin{figure*}[h]
\centering
\includegraphics[width=\textwidth]{Figures/exp5/fig_exp5_mechanism.pdf}
\caption{Source of \sysname{}'s performance gain. (a) Mean mission time divided into transit, communication dwell, and return/upload/dock time at (N=6) and the knee budget. (b) Time to the target accuracy for the same configurations. (c) RF classes selected by \sysname{}-Plan as the device population and mission budget change.}
\label{fig:mission-breakdown}
\end{figure*}

\section{Evaluation Results and Discussion}
\label{sec:results}

We compare \sysname{} with existing schedulers, examine where its gains come
from, study how it scales, and evaluate the contribution of its main design
choices. We then examine learning, channel variation, data heterogeneity,
payload size, and compute delay.

\subsection{Comparison with Existing Schedulers}
\label{sec:comparison}

Table~\ref{tab:scheduler_comparison} compares the evaluated approaches at
\(N=6\). With one mule and the knee budget, \sysname{}-Plan reaches the target
accuracy in 179~s, compared with 232--263~s for the baselines, a reduction of
23--32\%. FedEx-Merge takes 296~s and reaches the target in 92\% of trials.
\sysname{}-Adapt and \sysname{}-RL reach the target in 168 and 169~s,
respectively, but neither is significantly faster than \sysname{}-Plan. This
shows that most of the gain is already captured by the mission selected before
takeoff.

Increasing the budget to 300~s does not remove this difference.
\sysname{}-Plan still reaches the target in 179~s, while the baselines require
229--248~s. Under the tighter stress budget, the differences become smaller.
FedEx and FedEx-Merge remain significantly slower than \sysname{}-Plan, while
the other approaches differ by 37--64~s without statistical support. Most of
these approaches remain within the mission budget by serving fewer devices.
FedEx instead keeps a higher update yield by completing its full route and
therefore exceeds the budget in 84\% of missions.

\begin{table}[h]
\centering
\caption{Scheduler comparison at \(N=6\). Times are averaged over trials that
reach the target; parentheses report target reach when below 100\%.
For the one-mule stress setting, we also report mission-budget overrun and
updates merged per round. \(^{\dagger}\) Statistically significant difference
from \sysname{}-Plan. \(^{\S}\) Exploratory.}
\label{tab:scheduler_comparison}
\resizebox{\columnwidth}{!}{%
\begin{tabular}{l rrr rr rr}
\toprule
 & \multicolumn{3}{c}{$K=1$ mule} &
   \multicolumn{2}{c}{$K=3$ mules} &
   \multicolumn{2}{c}{$K=1$, stress} \\
\cmidrule(lr){2-4}\cmidrule(lr){5-6}\cmidrule(lr){7-8}
\textbf{Approach} & knee & stress & 300\,s & knee & stress & overrun & yield \\
\midrule
\sysname{}-Plan
& \textbf{179\,s} & \textbf{168\,s} & \textbf{179\,s}
& \textbf{53\,s} & \textbf{53\,s} & 29\% & 3.1 \\

\sysname{}-Adapt
& 168\,s & 165\,s & 168\,s
& 53\,s & 53\,s & 21\% & 3.1 \\

\sysname{}-RL$^{\S}$
& 169\,s & 166\,s & 168\,s
& 62\,s & 61\,s & 20\% & 3.1 \\
\midrule

Deadline-Heuristic
& 254\,s$^{\dagger}$ & 209\,s {\scriptsize(92\%)}
& 247\,s$^{\dagger}$ & 60\,s & 60\,s & 1\% & 2.0 \\

MAX-AoI
& 263\,s$^{\dagger}$ & 232\,s
& 247\,s$^{\dagger}$ & 60\,s & 60\,s & 2\% & 2.3 \\

Oort
& 235\,s$^{\dagger}$ & 236\,s {\scriptsize(95\%)}
& 238\,s$^{\dagger}$ & 64\,s & 64\,s & 4\% & 2.2 \\

Value-Aware FL
& 262\,s$^{\dagger}$ & 213\,s
& 244\,s$^{\dagger}$ & 63\,s & 63\,s & 2\% & 2.3 \\

FedEx
& 232\,s$^{\dagger}$ & 232\,s$^{\dagger}$
& 232\,s$^{\dagger}$ & 60\,s & 60\,s & 84\% & 3.4 \\

FedEx-Merge
& 296\,s {\scriptsize(92\%)}$^{\dagger}$
& 296\,s {\scriptsize(92\%)}$^{\dagger}$
& 296\,s {\scriptsize(92\%)}$^{\dagger}$
& 111\,s$^{\dagger}$ & 111\,s$^{\dagger}$ & 84\% & 3.4 \\

FedCS
& 236\,s$^{\dagger}$ & 206\,s {\scriptsize(98\%)}
& 229\,s$^{\dagger}$ & 60\,s & 60\,s & 2\% & 2.4 \\

Model-Aided RL
& 252\,s$^{\dagger}$ & 201\,s {\scriptsize(98\%)}
& 248\,s$^{\dagger}$ & 63\,s & 63\,s & 6\% & 2.6 \\
\bottomrule
\end{tabular}%
}
\end{table}

With three mules, most approaches reach the target within 53--64~s because each
mule serves a smaller region. FedEx-Merge remains slower at 111~s because its
merge rule applies older transported updates. The no-mule Live-FL reference
does not reach the target and ends at 0.693 accuracy, compared with 0.837 for
\sysname{}-Plan in the one-mule knee setting. The advantage of \sysname{} is
therefore strongest when a small number of mules must serve a larger
disconnected field.

\subsection{Source of the Performance Gain}
\label{sec:source-gain}

Figure~\ref{fig:mission-breakdown} shows that \sysname{} reaches the target
sooner mainly because it shortens each mission. At the \(N=6\) knee budget, the
approaches require a similar number of FL rounds, but the wide-band baselines
spend about 140~s traveling between contact locations. \sysname{}-Plan selects
the narrow class in 83\% of missions and reduces transit to about 18~s. The
additional communication time is smaller than the travel time saved, producing
missions that are 21--25\% shorter.

\textcolor{red}{Can we compare FedEx, and preferably Deadline-Heuristic, using narrow, medium, and/or the best fixed RF class? The main baselines currently use the wide class while FeRRy selects the RF class, and FeRRy pinned to wide is almost identical to FedEx. A reviewer may otherwise attribute the full gain to RF choice rather than joint RF and route planning.}

The fixed-band results separate the effect of RF choice from the remaining
planner decisions. When \sysname{}-Plan is restricted to the wide class, it
reaches the target in 232~s, the same as FedEx. With the narrow class, it reaches
the target in 184~s, close to the 179~s obtained when the planner selects the
class itself. The main gain in this setting therefore comes from choosing a
communication range that changes the route.

\textcolor{blue}{To separate RF choice from route planning, we also evaluated FedEx and
Deadline-Heuristic pinned to the medium and narrow classes, as an exploratory check
with 40 paired trials. At the \(N=6\) knee budget, this removes the gap: FedEx
reaches the target in 190 and 184~s and Deadline-Heuristic in 192 and 184~s,
compared with 190 and 184~s for \sysname{}-Plan pinned to the same classes. On the
narrow class, the three are identical in all 40 trials, because one stop serves
every device and leaves nothing to schedule, and \sysname{}-Plan's own choice of
class (179~s) is not significantly faster. In this setting, the gain therefore
comes from the RF class rather than from route planning. Planning matters more
under the stress budget. On the narrow class, FedEx exceeds the budget in 52\% of
missions, compared with 28\% for pinned \sysname{}-Plan, and Deadline-Heuristic
reaches the target in only 19 of 40 trials.}

The preferred class changes with the mission. As the number of devices grows or
the available mission time becomes smaller, \sysname{} moves from the narrow
class toward the medium and wide classes because communication time becomes
more important. There is therefore no single RF class that is best across all
settings.

\subsection{Scalability}
\label{sec:scalability}

\begin{figure*}[t]
\centering
\includegraphics[width=\textwidth]{Figures/exp5/fig_exp5_scale.pdf}
\caption{Scalability at each population's knee budget. Means are shown with
95\% bootstrap confidence intervals. The \(N=6\), one-mule setting uses
40 trials; the remaining settings use 20. (a) Time to target as the number of
devices increases with one mule. MAX-AoI, Oort, and FedCS are exploratory at
\(N=12\) and \(N=24\). (b) Planner execution time. (c) Time to target as the
number of mules increases. Hollow markers indicate that some trials did not
reach the target.}
\label{fig:scalability}
\end{figure*}

Figure~\ref{fig:scalability}a shows how performance changes as the number of
devices increases. At \(N=24\) and the knee budget, \sysname{}-Plan reaches the
target in 749~s, compared with 951~s for Deadline-Heuristic and 911~s for
Value-Aware FL, while FedEx requires 752~s. Under the stress budget,
\sysname{}-Plan, \sysname{}-Adapt, Deadline-Heuristic, and Value-Aware FL are
close at 651--662~s, and the difference between \sysname{}-Plan and FedEx is not
statistically supported. At \(N=12\), the evaluated approaches are also close.

The planning cost remains small as the population grows
(Figure~\ref{fig:scalability}b). Mean planning time increases from 0.04~s at
\(N=6\) to 1.53~s at \(N=96\), with a 95th-percentile time of 2.0~s at the
largest setting. In FerrySim, \sysname{} continues to serve about 54--59\% of
devices at the knee budget up to \(N=96\), showing that planner execution time
does not become the main bottleneck.

Adding mules provides another way to reduce collection time
(Figure~\ref{fig:scalability}c). At \(N=12\), \sysname{}-Plan reaches the target
in 486~s with one mule, 226~s with two, and 151~s with three. At \(N=24\), the
corresponding times are 749, 327, and 198~s. The improvement comes from dividing
the field into smaller regions that can be served at the same time.

The advantage over other schedulers becomes smaller as more mules are added.
Across the multi-mule scale settings, \sysname{}-Plan is usually faster on the
mean, but most differences are not statistically supported. At \(N=12\) with
three mules, \sysname{}-Adapt is significantly faster than \sysname{}-Plan,
reaching the target in 102 rather than 151~s. This suggests that decisions made
during flight become more useful when each mule works over a smaller region.

\subsection{Component Analysis}
\label{sec:components}

\begin{table*}[t]
\centering
\caption{Component analysis of \sysname{}. Each row removes or replaces one
design choice. Advantage is reported in the metric's original units and is
positive when \sysname{} performs better. Percentage differences are reported
in percentage points (pp).}
\label{tab:component-analysis}
\small
\begin{tabular}{
p{0.25\textwidth}
p{0.18\textwidth}
c
c
c
c
}
\hline
\textbf{Alternative} &
\textbf{Metric} &
\textbf{\sysname{}} &
\textbf{Alternative} &
\textbf{Advantage [95\% CI]} &
\textbf{Holm $p$} \\
\hline

\multicolumn{6}{c}{\textit{Coverage and participant freshness}} \\
\hline
Coverage term removed &
Round-close rate &
98\% & 73\% &
+25.0\,pp [+25.0, +25.0] &
$<0.001$ \\

Coverage term removed &
Age of updates &
0.78 & 1.10 &
+0.32 [+0.26, +0.38] &
0.002 \\

MAX-AoI &
Age of updates &
0.78 & 1.06 &
+0.29 [+0.16, +0.41] &
0.017 \\

Value-Aware FL &
Age of updates &
0.78 & 1.07 &
+0.30 [+0.18, +0.42] &
0.013 \\

Age cap removed &
Age of updates &
0.78 & 0.78 &
0 [0, 0] &
1.000 \\

\hline
\multicolumn{6}{c}{\textit{Mission objective and deadline handling}} \\
\hline

Dwell term removed &
Round-close rate &
98\% & 98\% &
0\,pp [0, 0] &
1.000 \\

FedCS-style round deadline &
Deadline-miss rate &
43.4\% & 48.7\% &
+5.3\,pp [+2.3, +8.5] &
0.010 \\

Oort-style preferred duration &
Deadline-miss rate &
43.4\% & 48.6\% &
+5.2\,pp [+2.1, +8.4] &
0.010 \\

\hline
\multicolumn{6}{c}{\textit{Aggregation and decisions during flight}} \\
\hline

FedBuff merge &
Time to target &
157\,s & 375\,s (26 seeds) &
+218\,s [+146, +295] &
0.004 \\

No per-stop rule (\(N=12\)) &
Time to target &
428\,s & 461\,s &
+33\,s [+8, +67] &
0.032 \\

\hline
\multicolumn{6}{c}{\textit{Backhaul selection}} \\
\hline

Fixed backhaul carrier &
Age of updates &
0.69 & 0.88 &
+0.19 [+0.13, +0.28] &
$<0.001$ \\

\hline
\end{tabular}

\vspace{2pt}
\footnotesize
Age of updates is measured in mule missions. The backhaul comparison uses the
time-varying backhaul setting; the remaining rows use the corresponding
settings described in Section~\ref{sec:environment}.
\end{table*}

Table~\ref{tab:component-analysis} shows that the coverage term has the clearest
effect on participant service. At \(N=12\), removing it reduces the round-close
rate from 98\% to 73\%. Under the stress budget, mean update age rises from
0.78 to 1.10 missions. MAX-AoI and Value-Aware FL reach 1.06 and 1.07 missions,
respectively. This shows that repeatedly missed devices need explicit priority
when serving them also requires UAV travel.

The deadline rule and decisions made during flight provide smaller but
measurable gains. Replacing the FeRRy deadline with FedCS- or Oort-style timing
raises the deadline-miss rate from about 43\% to 49\%. At \(N=12\), removing
the per-stop rule increases time to target from 428 to 461~s. Backhaul selection
also improves freshness when the backhaul changes over time, reducing mean
update age from 0.88 to 0.69 missions.
For studies using the 20-trial subset, the corresponding
\sysname{}-Plan mean at the \(N=6\) knee setting is 194~s, compared with
179~s over the 40 trials used in the main scheduler comparison.

Other components have little effect in the evaluated settings. Removing the age
cap or dwell term does not change their primary metrics, and several alternative
aggregation rules perform similarly because most collected updates are still
current. The results therefore place the main contribution at the mission level:
\textcolor{blue}{the plan's RF class provides} the largest performance gain, while coverage,
deadline handling, and decisions during flight improve participant service in
specific settings.

\subsection{Learning and Decisions During Flight}
\label{sec:learning}

\begin{figure*}[t]
\centering
\includegraphics[width=\textwidth]{Figures/exp5/fig_exp5_learning.pdf}
\caption{Effect of learning inside the FeRRy mission plan.
(a) Time to target for \sysname{}-Plan, \sysname{}-Adapt, \sysname{}-RL,
and Model-Aided RL under the knee and stress budgets.
(b) Mean mission time at the knee budget divided into transit,
communication, and return/upload/dock time.}
\label{fig:learning-comparison}
\end{figure*}

Figure~\ref{fig:learning-comparison} compares the planned, deterministic, and
learned FeRRy variants with Model-Aided RL. \sysname{}-RL does not outperform
\sysname{}-Adapt. At \(N=6\), \sysname{}-Plan also performs similarly, showing
that most of the gain is determined by the mission selected before takeoff.

The comparison with Model-Aided RL further separates planning from learning.
At the \(N=6\) knee budget, \sysname{}-RL reaches the target in 169~s, compared
with 252~s for Model-Aided RL. The difference comes mainly from travel:
\sysname{}-RL inherits the RF choice from the mission plan and spends about
17~s in transit, while Model-Aided RL uses the wide class and spends about
142~s in transit. Under the stress budget, the difference is not statistically
significant.

At \(N=12\), we also vary the channel period around the pace of the decisions.
\sysname{}-Adapt remains significantly faster at the knee budget by 78--89~s,
while the differences under stress are not statistically supported. These
results do not show an advantage from adding learning to the FeRRy plan. For the
conditions we evaluate, the deterministic rule is sufficient.

\textcolor{blue}{Changing band also changes propulsion energy, because it changes how long the mule
hovers at a stop. In an exploratory comparison paired by seed, \sysname{}-Adapt
uses 11\% less energy per mission than \sysname{}-Plan at \(N=6\) and the knee
budget (20.1 against 22.6~kJ), 10\% less under stress, and 5\% less at \(N=12\) and
the knee budget, but 5\% more at \(N=24\) and the knee budget (48.1 against
45.6~kJ). The differences at \(N=12\) and \(N=24\) under stress are not
statistically supported. Moves to a wider class shorten the stop, while moves from
medium to narrow lengthen it. The rule compares classes using the SNR read on
arrival, but each device is then served at its later SNR: among first-mission
changes from medium to narrow, narrow looked faster at all 25 stops and was faster
at 5. Uplink losses in our model do not depend on the RF class, so these savings
may be optimistic.}

\subsection{Robustness and Limitations}
\label{sec:robustness}

\textcolor{red}{What exactly does ``17 dB interference'' mean in this experiment? Which channel-model parameter is changed, and by how much? Please give the precise setting.}

\begin{table*}[t]
\centering
\caption{Robustness at \(N=6\) and the knee budget. Times are averaged over
trials that reach the target; parentheses report target reach when below 100\%.
A shorter time with a low reach rate should therefore not be interpreted as
better overall performance. \(^{\S}\) Exploratory.}
\label{tab:robustness}
\small
\setlength{\tabcolsep}{5pt}

\begin{tabular}{llcccc}
\hline
\textbf{Category} &
\textbf{Condition} &
\textbf{\sysname{}-Plan} &
\textbf{\sysname{}-Adapt} &
\textbf{Deadline-Heuristic} &
\textbf{Additional Reference} \\
\hline

\multirow{2}{*}{Compute}
& Training spread\textcolor{blue}{$^{\S}$}
& \textcolor{blue}{209}
& \textcolor{blue}{204}
& \textcolor{blue}{276}
& \textcolor{blue}{FedCS: 276; FedEx: 298} \\

& 20\% stragglers\textcolor{blue}{$^{\S}$}
& \textcolor{blue}{199}
& \textcolor{blue}{193}
& \textcolor{blue}{289}
& \textcolor{blue}{FedCS: 270; FedEx: 309} \\
\hline

Data
& Dirichlet \(\alpha=0.1\)
& 378 (55\%)
& 329 (55\%)
& 507 (55\%)
& Value-Aware FL: 519 (50\%) \\
\hline

\multirow{3}{*}{Channel}
& \textcolor{blue}{\(A=17\) dB}
& 255 (95\%)
& 209 (95\%)
& 332
& FedEx: 276$^{\S}$; FedCS: 315 (95\%)$^{\S}$ \\

& \(n=2.7\)
& 257
& 263
& 261
& FedEx: 237 (90\%)$^{\S}$; FedCS: 296 (95\%)$^{\S}$ \\

& \(\sigma=8\) dB
& 219 (90\%)
& 183 (85\%)
& 281
& FedEx: 262$^{\S}$; FedCS: 292 (95\%)$^{\S}$ \\
\hline
\end{tabular}

\vspace{2pt}
\footnotesize
The channel experiments use the time-varying backhaul setting.
\textcolor{blue}{The compute experiments run for 10 missions with 40 paired trials.}
When device training time is modeled, \sysname{}-Plan merges \textcolor{blue}{0.99 updates
per round with training spread and 0.89 with stragglers}, compared with 3.54 when
training delay is not included.
\end{table*}

Table~\ref{tab:robustness} shows that data heterogeneity makes the learning task
harder for every approach. With strongly non-IID data
(\(\alpha=0.1\)), only about half of the trials reach the target.
Among those trials, \sysname{}-Plan reaches it in 378~s and
\sysname{}-Adapt in 329~s, compared with 507~s for Deadline-Heuristic and
519~s for Value-Aware FL. At \(\alpha=1\), the reach rate improves to
65--75\%. With IID data, all approaches reach the target.

\textcolor{red}{Please define the compute experiments. What distribution/range is used for ``training spread,'' and how are the 20\% stragglers generated? We need enough detail to reproduce Table~\ref{tab:robustness}.}

\textcolor{blue}{In the interference condition, we raise the amplitude $A$ in
Eq.~\eqref{eq:channel-model} from 5 to 17~dB on the contact link, keeping
$\sigma_I=1.5$~dB and $P_c=60$~s. A pilot selected 17~dB as the harshest tested
level at which Deadline-Heuristic still served at least half of its default share
of devices (41.7\% against 57.7\%).}
The channel results show that the advantage is not universal.
\sysname{}-Adapt often improves on the original plan by reacting to the
observed channel, but under \(n=2.7\), \sysname{}-Plan,
\sysname{}-Adapt, and Deadline-Heuristic require similar time. FedEx reaches
237~s in this setting compared with 257~s for \sysname{}-Plan, although the
difference is not statistically supported. These results show that the benefit
depends on how much the communication range changes the route.

Payload size has a similar effect. With the measured 18.8~KB model,
\sysname{}-Plan uses the narrow class and reaches the target in 59~s, compared
with 272~s for Deadline-Heuristic, 239~s for FedEx, and 243~s for FedCS.
At 1~MB, the narrow class is still selected for most contacts and FeRRy keeps a
clear advantage. As the payload grows, communication time becomes more
important and the planner shifts toward the medium and wide classes. At 10~MB,
it uses only the wide class and reaches the target in 344~s. No baseline differs
significantly from \sysname{}-Plan at 2, 5, or 10~MB. The benefit of RF-based
route planning therefore becomes smaller as the model update grows.

\textcolor{blue}{Compute delay reduces how many updates each mission collects, but \sysname{}
keeps its lead. In the compute experiments, local training takes simulated time
$T_j=68\,e^{0.5z_j}$~s, where $z_j$ is a standard normal draw fixed per device and
trial, so all methods see the same times. Training restarts whenever a device
receives a model, and the first run starts at the first takeoff. A contact before
training finishes collects no update. A pilot selected the 68~s median, at which
21\% of contacts after the first mission found no update ready. The straggler
setting multiplies these times by 5 for one randomly chosen device of six.
The other rows of Table~\ref{tab:robustness} end after 4 missions, which here gives
the methods unequal mission time: \sysname{}-Plan flies 312--335~s against
572--853~s for the baselines and reaches the target in 70--75\% of trials against
85--100\%. We therefore report these two settings over 10 missions, as an
exploratory comparison with 40 paired trials. Every method then reaches the target
in every trial. \sysname{}-Plan needs 209~s with distributed training times and
199~s with stragglers, against 276--317~s and 270--334~s for the baselines. It is
significantly faster than Value-Aware FL in both settings, than FedCS with
distributed training times, and than FedEx with stragglers, and no baseline is
faster. Its updates per round fall from 3.54 to 0.99 and 0.89 because the mule may
arrive before a device has finished training, but its missions average 90--93~s
against 173--216~s, so it merges updates at about the same rate per second of
flight. These results show that \sysname{} tolerates training delay, not that it
handles it better than the baselines: it already leads them by 53--83~s without
training delay (Table~\ref{tab:scheduler_comparison}), and FedCS, which selects
devices using their training times, is not faster.}
\sysname{} currently plans using communication and mobility information but
does not consider when each local update will be ready. Including expected
training completion time in the planner is therefore an important extension.

Finally, the mission plan uses the expected channel before takeoff. Under the
stress budget, this leads to budget overruns in 29\% of missions, with an
average overrun of 33~s and a median of 20~s. Reserving additional mission time
or accounting for channel variation during planning could reduce these
overruns. UAV movement, wireless conditions, and propulsion energy are also
modeled rather than measured. Validation with measured UAV channels and battery
limits remains necessary before drawing conclusions about physical deployment.

% \subsection{Against the State of the Art}
% \label{sec:res:sota}

% With one mule at the knee budget, F reaches $\tau$ in 194\,s and each of the
% seven baselines needs 264--282\,s (Table~\ref{tab:exp5_headline}): F is
% 27--31\% faster, every comparison is a claim, and F gets there first in
% 80--90\% of paired seeds. FedEx with its own merge (D4$_{\text{FedEx}}$) is as
% slow (281\,s) but misses $\tau$ once, so its comparison is not a claim. FX
% (178\,s) and FQ (180\,s) do not differ from F; FQ, flown later, also beats D4
% and E3 at the knee in its exploratory family. Under the stress budget only the
% FedEx route is significantly slower: the other baselines trail by 28--60\,s but
% stay within the budget by leaving devices out. With three mules every scheduler
% reaches $\tau$ in 52--65\,s except D4$_{\text{FedEx}}$ (101\,s), whose
% unweighted, one-tour-old updates slow learning. Without a mule (H0), FL over the
% degraded link merges 0.5 updates per round and ends at 0.686 accuracy, barely
% above the all-benign 0.667, against F's 3.5 and 0.831.


% \subsection{Why FeRRy Is Faster: Reach}
% \label{sec:res:mechanism}

% F does not need fewer rounds (every arm reaches $\tau$ in 1.65--1.84 rounds,
% with near-identical update yield); its rounds are shorter
% (Fig.~\ref{fig:exp5_mechanism}). Every baseline flies the wide band, which
% reaches a device only from nearby, and spends about 140\,s of a mission in
% transit. F's plan chooses the narrow band in 89\% of missions, whose range
% covers the field from about one stop: F spends 15\,s in transit and 107\,s
% dwelling, and its missions are 22--26\% shorter. FX and FQ inherit this reach
% (narrow or medium on 95\% and 98\% of missions; Fig.~\ref{fig:exp5_learning}b).
% Pinned to the wide band, F reaches $\tau$ in 265\,s, D4's own time; pinned to
% narrow, in 196\,s. (These times are descriptive: the band-pinning study's
% pre-registered metric, the age of updates, does not differ.) The choice is near
% the best, serving as many devices as an exhaustive oracle at the knee budget and
% 4 points fewer under stress, and it adapts (Fig.~\ref{fig:exp5_mechanism}c): as the
% field fills and the budget tightens, F shifts to medium and wide, mostly wide at
% $N=24$ under stress.


% \subsection{The Design Claims}
% \label{sec:res:claims}

% At $N=24$ (Table~\ref{tab:exp5_scale}), F is significantly faster than H1 and D3
% at the knee (749 against 951 and 911\,s), ties the FedEx route (752\,s), and
% beats it under the stress budget (661 against 752\,s); at $N=12$ no comparison
% is a claim, and FQ (461 and 389\,s) does not differ from FX (424 and 386\,s).
% Planning stays cheap (Table~\ref{tab:exp5_cost}): alone on one host, a plan
% takes 0.04\,s at $N=6$ and 1.5\,s at $N=96$ (p95 2.0\,s). In FerrySim the plan
% serves 54--57\% of the devices at each scale's knee budget up to $N=96$, and at
% $N=12$ more mules cut F's time to $\tau$ from 486\,s to 226\,s with two and
% 151\,s with three (Table~\ref{tab:exp5_mules}).




% Five tests hold (Table~\ref{tab:exp5_claims}). Without the coverage term the
% round-close rate falls from 99\% to 74\% (C2). Under the stress budget FeRRy's
% age of updates is 0.78 missions, against 1.06 for MAX-AoI, 1.07 for the Whittle
% index and 1.10 without the coverage term (C5). On the time-varying backhaul the
% adaptive controller lowers it from 0.88 to 0.69, close to the most adaptation
% can give, since the controller sees the SNR that decides each loss. Other parts
% show no measurable effect here: pinning the band leaves the age of updates
% unchanged (its effect is on time), the dwell term and the age cap change nothing
% over four missions, and the merge rules tie because collected updates almost
% always have merge age zero (FedBuff's buffering slows learning, and FedProx is
% not significantly faster). The selection literature's deadline forms miss
% slightly more deadlines (47\% against 42\% under stress), but not significantly.

% \begin{table*}[t]
% \centering
% \caption{Tests of the design claims, each against an alternative that replaces one part of FeRRy's design. Advantage: the paired difference as a share of FeRRy's value, positive when FeRRy's design does better, with its 95\% CI. \emph{Holds}: significant at Holm $p<0.05$.}
% \label{tab:exp5_claims}
% \footnotesize
% \resizebox{\textwidth}{!}{%
% \begin{tabular}{l l rr r r l}
% \toprule
% \textbf{Alternative tested} & \textbf{Metric} & \textbf{FeRRy} & \textbf{Alternative} & \textbf{FeRRy's advantage [95\% CI]} & \textbf{Holm $p$} & \textbf{Verdict} \\
% \midrule
% \multicolumn{7}{l}{\textit{C1: reach is a decision}} \\
% \quad Band class pinned to wide & age of updates (missions) & 0.70 & 0.72 & +3\% [0, +7] & 1.000 & no difference \\
% \quad Band class pinned to narrow & age of updates (missions) & 0.70 & 0.70 & 0\% [0, 0] & 1.000 & no difference \\
% \midrule
% \multicolumn{7}{l}{\textit{C2: one derived objective}} \\
% \quad Coverage term removed ($N{=}12$) & round-close rate & 99\% & 74\% & +25\% [+25, +25] & $<$0.001 & \textbf{holds} \\
% \quad Dwell term removed ($N{=}12$) & round-close rate & 99\% & 99\% & 0\% [0, 0] & 1.000 & no difference \\
% \midrule
% \multicolumn{7}{l}{\textit{C3: one deadline, three roles}} \\
% \quad Plain-mean merge & time to $\tau$ & 194\,s & 194\,s & 0\% [0, 0] & 1.000 & no difference \\
% \quad Async-HFL merge & time to $\tau$ & 194\,s & 194\,s & 0\% [0, 0] & 1.000 & no difference \\
% \quad FedBuff merge & time to $\tau$ & 159\,s & 346\,s (13 seeds) & +117\% [+59, +183] & 0.117 & no difference \\
% \quad FeRRy's merge with FedProx & time to $\tau$ & 194\,s & 146\,s & $-$25\% [$-$64, +10] & 1.000 & no difference \\
% \quad Round deadline, after FedCS (stress) & deadline-miss rate & 42.2\% & 47.4\% & +13\% [+3, +23] & 0.152 & no difference \\
% \quad Preferred duration, after Oort (stress) & deadline-miss rate & 42.2\% & 47.2\% & +12\% [+3, +22] & 0.152 & no difference \\
% \midrule
% \multicolumn{7}{l}{\textit{C4: two clocks, re-decided per stop}} \\
% \quad Learned per-stop score, FQ ($N{=}12$) & time to $\tau$ & 424\,s & 461\,s & +9\% [0, +23] & 0.555 & no difference \\
% \midrule
% \multicolumn{7}{l}{\textit{C5: fairness under physical cost}} \\
% \quad MAX-AoI scheduler, D1 (stress) & age of updates (missions) & 0.78 & 1.06 & +37\% [+21, +53] & 0.009 & \textbf{holds} \\
% \quad Whittle index, D3 (stress) & age of updates (missions) & 0.78 & 1.07 & +38\% [+23, +55] & 0.007 & \textbf{holds} \\
% \quad Coverage term removed (stress) & age of updates (missions) & 0.78 & 1.10 & +41\% [+34, +49] & 0.001 & \textbf{holds} \\
% \quad Age cap removed (stress) & age of updates (missions) & 0.78 & 0.78 & 0\% [0, 0] & 1.000 & no difference \\
% \midrule
% \multicolumn{7}{l}{\textit{Component ablations}} \\
% \quad Whole stops only (stress) & age of updates (missions) & 0.81 & 0.90 & +11\% [+4, +18] & 0.078 & no difference \\
% \quad Fixed backhaul carrier, no L1 & age of updates (missions) & 0.69 & 0.88 & +28\% [+18, +40] & $<$0.001 & \textbf{holds} \\
% \bottomrule
% \end{tabular}%
% }
% \end{table*}
% \subsection{Learning Inside the Plan}
% \label{sec:res:rl}

% The learned-score study is flat: no $\gamma$ beats $\gamma=0$ by the
% pre-registered margin, and the best ($\gamma=0.75$) matches FX while trailing a
% one-step greedy rule. In the full system FQ is again indistinguishable from FX
% (C4) and, at $N=6$, from F, so FX is FeRRy's per-stop rule. The two decide on
% different information (FQ before serving, from the priced dwell; FX after, from
% the realized mission time), so this compares two rules, not two rankings.

% The exploratory comparison sets FQ against E3, a single-agent DQN after Chen et
% al.~\cite{chen2023model} that picks every next stop with no plan, on the same 20
% seeds (Fig.~\ref{fig:exp5_learning}). At the knee FQ reaches $\tau$ in 180\,s
% against E3's 280\,s (difference 100\,s, 95\% CI [63, 141], Holm $p<0.001$ across
% the eight exploratory comparisons); under stress it is faster over 19 shared
% seeds (174 against 211\,s), but the CI includes zero. FQ inherits the plan's
% band and spends 15\,s of a 155\,s mission in transit, against E3's 144\,s of
% 228\,s, and in FerrySim E3's held-out return falls 6\% below FX's while FQ's
% matches it. The gain comes from the plan, not from learning, which supports the
% two-clock structure (C4).


% \subsection{Robustness and Limitations}
% \label{sec:res:robust}

% With strongly non-IID data (Dirichlet $\alpha=0.1$;
% Table~\ref{tab:exp5_robust}), every arm reaches $\tau$ in about half the trials,
% and F stays significantly faster than H1 and D3 where both do. The radio study
% flies the time-varying backhaul, so its default (F at 216\,s) is not the main
% configuration (194\,s); under its harsher channels F stays ahead of H1 without a
% claim, and FX is significantly faster than F on the default channel and with
% 8\,dB shadowing. Two limitations stand out. FeRRy's plan prices the mean SNR, so
% under the stress budget 30\% of its missions overrun, by 28\,s on average
% (median 20\,s; Table~\ref{tab:exp5_headline}); the FedEx route, which never skips
% a device, overruns in 85\%, by 22\,s on average, and F keeps nearly its update
% yield with fewer stops. And the
% plan does not price training time: when devices take time to train, F's yield
% falls from 3.5 to 0.7 updates per round and it reaches $\tau$ in only 70--75\% of
% trials, while H1, D4 and D5 keep 85--100\%.

% \begin{table*}[t]
% \centering
% \caption{Robustness at $N=6$ and the knee budget: time to $\tau$. Notation as in Table~\ref{tab:exp5_headline}; $^{\ast}$significantly faster than F.}
% \label{tab:exp5_robust}
% \small
% \begin{tabular}{l rr r rrrr}
% \toprule
%  & \multicolumn{2}{c}{compute (Study 10)} & data (Study 11) & \multicolumn{4}{c}{channel (Study 13)} \\
% \cmidrule(lr){2-3}\cmidrule(lr){4-4}\cmidrule(lr){5-8}
% \textbf{Arm} & train spread & stragglers & Dir(0.1) & default & 17 dB interf. & $n_{\text{pl}}{=}2.7$ & $\sigma{=}8$ dB \\
% \midrule
% F & \textbf{137\,s}\,{\scriptsize(75\%)} & \textbf{159\,s}\,{\scriptsize(70\%)} & \textbf{378\,s}\,{\scriptsize(55\%)} & \textbf{216\,s}\,{\scriptsize(95\%)} & \textbf{255\,s}\,{\scriptsize(95\%)} & \textbf{257\,s} & \textbf{219\,s}\,{\scriptsize(90\%)} \\
% FX & 120\,s\,{\scriptsize(75\%)} & 146\,s\,{\scriptsize(70\%)} & 329\,s\,{\scriptsize(55\%)}$^{\ast}$ & 192\,s\,{\scriptsize(95\%)}$^{\ast}$ & 209\,s\,{\scriptsize(95\%)} & 263\,s & 183\,s\,{\scriptsize(85\%)}$^{\ast}$ \\
% H1 & 325\,s$^{\dagger}$ & 296\,s\,{\scriptsize(95\%)} & 507\,s\,{\scriptsize(55\%)}$^{\dagger}$ & 274\,s & 332\,s & 261\,s & 281\,s \\
% D3 (Cui) & -- & -- & 519\,s\,{\scriptsize(50\%)}$^{\dagger}$ & -- & -- & -- & -- \\
% D4 (FedEx route) & 333\,s$^{\dagger}$ & 304\,s\,{\scriptsize(90\%)} & -- & -- & -- & -- & -- \\
% D5 (FedCS) & 222\,s\,{\scriptsize(95\%)}$^{\dagger}$ & 250\,s\,{\scriptsize(85\%)} & -- & -- & -- & -- & -- \\
% \bottomrule
% \end{tabular}
% \end{table*}

\section{Conclusion and Future Work}
\label{sec:conclusion}

This paper presented \sysname{}, a framework for federated learning in
disconnected environments using UAVs as data mules. \sysname{} selects the
communication band and UAV route before each mission based on device deadlines,
communication time, and the mission budget. During flight, it can change the
band or revise the remaining route when the observed channel differs from the
planned condition. The UAV collects local model updates, carries them back to
the edge server, and later delivers the updated global model to the devices.

We evaluated \sysname{} using a multi-process prototype that trains a real
intrusion detection model on CICIoT2023, while UAV movement, wireless
conditions, mission time, and energy were simulated. With six devices, one UAV,
and a 1~MB model update, \sysname{} reached the target accuracy 23--32\% sooner
than the compared schedulers. The main gain came from selecting a narrow band
with greater range, which reduced UAV travel between devices. \textcolor{blue}{Those schedulers used the wide band; pinned to the narrow band, FedEx and
Deadline-Heuristic were not significantly slower.} Under a tight
mission budget, \sysname{} also kept updates fresher and stayed within the
budget in 71\% of missions; the remaining missions exceeded it by 33~s on
average. The advantage became smaller as model updates grew or more UAVs were
added. At 10~MB, \sysname{} selected the wide band and its lead was no longer
statistically significant, while with two or three UAVs most approaches
performed similarly.

The current planner does not account for device training time or uncertainty in
the channel before takeoff. Future work will study these factors directly and
evaluate \sysname{} with measured UAV wireless links and battery limits in a
physical testbed.

% IPDPS 2027 asks authors to disclose AI-generated text in the Acknowledgements.
% DRAFT wording: the authors must confirm or rewrite it. Funding acknowledgements
% stay out until the camera-ready (double-anonymous review).
\section*{Acknowledgment}
Parts of this paper's text, and of the code for its experiments and analysis,
were drafted with the assistance of an AI model (Claude, Anthropic). The authors
reviewed, edited, and verified all of it, and take responsibility for the
content.

\bibliographystyle{ieeetr}
\bibliography{biblio}

\end{document}
